\chapter{多模态数据的对齐与融合}
\label{chap:vision-multimodal-alignment}

设一段10秒的街口录像同时保存了三样东西：每秒25帧的画面，每秒16000个采样的录音，以及自动生成的一行字幕“红色汽车右转后鸣笛”。我们想回答：鸣笛的是不是那辆右转的红车？画面显示一辆红车在第3至第6秒转弯，录音在第7秒附近有一声尖锐的鸣响，字幕的词没有时间戳，也不说明“鸣笛”属于哪辆车。回答这个问题至少要做三件事。第一，确定三种记录中的哪些单元说的是同一件事：哪几帧对应“红色汽车”，哪一段录音对应“鸣笛”，鸣响时刻对应哪几帧。第二，把对上的证据合成一个判断：声音说有鸣笛，画面说红车刚转完弯，两者共同支持还是相互矛盾。第三，处理不完整的情形：录音被风声淹没，字幕缺失，或者训练时模型学会只看画面而忽略声音。

本章把这三件事分别称为对齐、融合和缺失与不平衡下的学习。第一节先说明多模态数据为什么难以直接比较：各模态的数值形式和时间分辨率不同，配对关系往往只在整段层面成立，而且某些模态可能缺席或彼此矛盾。这些特点决定了后面三节要回答的问题。第二节讨论怎样建立跨模态对应：典型相关分析和对比学习回答“怎样把两种表示放到可以比较的空间”，动态时间规整和注意力回答“在可比较之后，哪个单元对应哪个单元”。第三节在对应已经建立的前提下讨论怎样合并证据，三种融合位置在参数量与交互能力之间作出不同取舍。第四节放松“所有模态都在且都被充分利用”的假设，讨论训练和推理中模态缺失、以及优化中模态贡献失衡时的做法。

阅读本章需要四项前置知识，用到时都会先写出所需结论：卷四降维章中的典型相关分析，卷一图论章中的动态规划最短路，卷五的注意力机制，卷二概率分布章中的互信息与对比损失下界。本章讲的是对齐与融合的原理，具体模型放在后面：CLIP与ImageBind的结构见\ref{chap:vision-multimodal-models}章，跨模态检索与定位见\ref{chap:vision-multimodal-understanding}章，音画同步生成见\ref{chap:vision-multimodal-generation}章。文中的小向量、概率和代价表均为教学构造，不是任何模型的测量结果。

\section{多模态数据的特点}
\label{sec:vision-alignment-characteristics}

开篇的问题之所以不能直接交给一个网络，是因为三种记录在三个方面不对等。它们的数值形式和单元粒度不同，不能逐个相减；它们之间只知道“同属一段录像”，不知道哪一帧对应哪个词；某种记录还可能缺失或说了与其他记录相反的话。本节逐一说明这三个特点，并指出它们各自对应本章哪一节要解决的问题。

Baltru\v{s}aitis等人的综述把多模态学习的核心挑战归纳为五类：表示、转换、对齐、融合与协同学习\scite{vision-mmsurvey2019}表示指怎样把各模态编码成网络可用的数值；转换指从一种模态生成另一种模态，例如由图像生成描述；对齐指找出不同模态单元之间的对应；融合指合并多种模态的信息来预测；协同学习指借助一种模态的信息帮助另一种模态的学习，例如某种模态标注稀少时。本卷把表示放在\ref{chap:vision-image-video-models}章和\ref{chap:vision-multimodal-models}章，把转换放在多模态理解与生成两章。本章讨论对齐与融合，并在最后一节讨论缺失与不平衡，它与协同学习关心的“一种模态怎样帮助或拖累另一种模态”有关。这里借用综述的划分来确定本章边界，具体组织仍按本书的问题顺序展开。

\subsection{表示的异构}
\label{sec:vision-alignment-heterogeneity}

\term{模态}（modality）指信息的类型或通道。同一事件的不同模态，记录成数值后在三个方面不同：取值是连续的还是离散的，一个计算单元覆盖多大的空间或时间范围，相邻单元之间的距离有什么含义。我们先把开篇录像的三种记录写成张量，再看这些差别具体落在哪里。

画面按帧记录。帧率为$f_{\mathrm v}$、时长为$T$秒、每帧$h\times w$个像素的三通道视频是一个四阶数组$\bm V\in\{0,\ldots,255\}^{n_{\mathrm v}\times h\times w\times 3}$，其中$n_{\mathrm v}=f_{\mathrm v}T$。像素值是有序的整数，相邻像素的值通常接近，数值差能反映亮度差。录音按采样点记录，采样率为$f_{\mathrm s}$时，波形是长度$f_{\mathrm s}T$的一维序列$\bm s$，相邻采样相隔$1/f_{\mathrm s}$秒；单个采样值几乎不携带语义，要几百个采样连起来才构成可辨的声音。字幕是离散的\term{词元}（token）序列$(a_1,\ldots,a_m)$，$a_j$是词表中的编号；两个编号之差没有意义，“汽车”的编号比“自行车”大1，不说明它们更相近。

把声音交给网络之前，常先算成\term{梅尔频谱}（mel-spectrogram）：用短时窗在波形上滑动，每个窗算一次频谱，再按接近人耳听觉的梅尔刻度把频率分成若干组，记录每组能量。窗长25毫秒、窗间隔10毫秒时，每秒得到100个梅尔帧。表\ref{tab:vision-alignment-shapes}按这些常见设置列出开篇10秒录像的张量形状。字幕的词元数取决于分词器，表中按“红色／汽车／右转／后／鸣笛”计为5个，是教学设定。

\begin{table}[htbp]
  \centering
  \small
  \begin{tabularx}{\linewidth}{@{}p{18mm}p{30mm}p{26mm}X@{}}
    \toprule
    记录 & 10秒的张量形状 & 单元时长 & 取值与相邻单元关系\\
    \midrule
    视频帧 & $250\times224\times224\times3$ & 40毫秒 & 8位整数；相邻像素、相邻帧的值通常接近\\
    波形 & $160000$ & 0.0625毫秒 & 实数；单个采样几乎不含语义\\
    梅尔频谱 & $1000\times128$ & 10毫秒（窗间隔） & 非负能量；纵轴是频率而非空间\\
    字幕词元 & $5$（教学设定） & 不固定 & 离散编号；编号之差无意义\\
    \bottomrule
  \end{tabularx}
  \caption{开篇10秒街口录像的各模态张量。视频按25帧每秒、缩放到$224\times224$，波形按16kHz，梅尔频谱按10毫秒窗间隔与128个频率组；这些是常见设置，不是唯一约定。}
  \label{tab:vision-alignment-shapes}
\end{table}

单元时长的差别直接造成对应上的困难。图\ref{fig:vision-alignment-time-resolution}把开篇录像中的一秒放大：这一秒有25个视频帧、100个梅尔帧和16000个波形采样，而整段10秒的字幕一共只有5个词元。假设字幕来自同步的旁白，图中画出其中三个词元在这一秒内被说出时各自覆盖的时间，一个词元覆盖许多帧，各词元的时长也不同。若要判断“鸣笛”这个词对应哪些帧，就需要知道它覆盖的起止时刻，而自动字幕往往只给出整句或整段的时间范围。

\begin{figure*}[tbp]
  \centering
  % 一秒教学时间轴；宽约125mm。帧率、窗间隔为常见设置，词元时刻为构造值。
\begin{tikzpicture}[x=11cm,y=1cm,font=\figurefont\small]
  % 视频帧：25帧/秒，每帧间隔40毫秒
  \node[anchor=east] at (-0.02,3.0) {视频帧};
  \foreach \k in {0,...,24}{
    \fill[FigureInput!25,draw=FigureInput,line width=0.4pt]
      ({\k*0.04+0.004},2.78) rectangle ({\k*0.04+0.036},3.22);
  }
  \node[anchor=west,text=BookMuted] at (1.02,3.0) {25个，间隔40毫秒};
  % 梅尔帧：窗间隔10毫秒
  \node[anchor=east] at (-0.02,2.0) {梅尔帧};
  \foreach \k in {0,...,99}{
    \draw[FigureModel,line width=0.5pt] ({\k*0.01+0.005},1.8) -- ({\k*0.01+0.005},2.2);
  }
  \node[anchor=west,text=BookMuted] at (1.02,2.0) {100个，间隔10毫秒};
  % 波形采样：16kHz，过密而以色带表示
  \node[anchor=east] at (-0.02,1.0) {波形采样};
  \fill[FigureOutput!20] (0,0.85) rectangle (1,1.15);
  \node[anchor=west,text=BookMuted] at (1.02,1.0) {16000个};
  % 文字词元：构造的起止时刻
  \node[anchor=east] at (-0.02,0.0) {字幕词元};
  \foreach \a/\b/\w in {0.05/0.38/右转,0.38/0.52/后,0.62/0.95/鸣笛}{
    \draw[BookTeal,fill=BookTeal!10,line width=0.6pt,rounded corners=1.8pt]
      (\a,-0.25) rectangle (\b,0.25);
    \node at ({(\a+\b)/2},0) {\w};
  }
  \node[anchor=west,text=BookMuted] at (1.02,0.0) {本秒3个，长短不一};
  % 时间轴
  \draw[-{Stealth[length=2mm]},BookInk,line width=0.8pt] (0,-0.65) -- (1.05,-0.65);
  \foreach \t/\lab in {0/0,0.25/0.25,0.5/0.5,0.75/0.75,1/1}{
    \draw[BookInk,line width=0.7pt] (\t,-0.6) -- (\t,-0.7);
    \node[below] at (\t,-0.7) {\lab};
  }
  \node[anchor=north east] at (1.05,-1.0) {本秒内的时间（秒）};
\end{tikzpicture}

  \caption{同一秒内各模态的单元。视频帧与梅尔帧按常见的固定间隔排列，波形采样过密而以色带表示；字幕词元假设来自同步旁白，其起止时刻为教学构造，用来说明一个词可以覆盖许多帧，且各词时长不同。读图时比较各行单元的宽窄。}
  \label{fig:vision-alignment-time-resolution}
\end{figure*}

固定帧率的模态之间，时间对应可以靠重采样建立。视频第$k$帧的时刻是$(k-1)/f_{\mathrm v}$秒，梅尔第$l$帧的中心约为$(l-1)\Delta_{\mathrm h}+\Delta_{\mathrm w}/2$秒，其中$\Delta_{\mathrm h}$是窗间隔、$\Delta_{\mathrm w}$是窗长。以25帧每秒、10毫秒间隔为例，每个视频帧对应4个梅尔帧；把4个梅尔帧的特征取平均或拼接，就得到与视频帧同频的音频特征。这一换算只在两条记录的时钟同步时成立。录制设备之间常有几十毫秒的固定偏移，视频还可能丢帧或变帧率，此时需要用真实时间戳，或者用下一节的同步方法估计偏移。帧率与时间采样的基础见\ref{sec:vision-videofound-frame-rate}节。

另一类差别不能靠换算消除。像素网格上两个位置的欧氏距离近似反映空间远近，梅尔频谱上两个格子却可能一个在时间轴、一个在频率轴上相邻；文字词元之间没有任何坐标距离。因此各模态需要各自的编码器，把原始记录变成维数相同的向量后才能计算相似度。即便如此，两个编码器输出的第$k$维也没有共同含义：图像编码器的第5维和文字编码器的第5维是各自训练形成的坐标。让这些坐标在语义上对应，正是\ref{sec:vision-alignment-alignment}节中典型相关分析与对比学习要做的事。

\subsection{配对的弱与噪声}
\label{sec:vision-alignment-weak-pairing}

训练对齐需要知道哪些内容相互对应。人工标出“第3至第6秒的红车对应字幕中的‘红色汽车’”成本很高，大规模多模态数据通常只给出更粗的配对：这张网页图片与这段替代文字出现在一起，这段视频与这段旁白在同一时间段播放。我们把这种只在整体层面成立、内部单元对应未知的配对称为\term{弱配对}（weak pairing）。

可以用一个对齐矩阵把“弱”写清楚。设视觉一侧有$n$个单元，例如帧或区域，文字一侧有$m$个单元，例如词。若知道每对单元是否对应，可以记$\bm A\in\{0,1\}^{n\times m}$，$A_{i,j}=1$表示第$i$个视觉单元与第$j$个词对应。弱配对只告诉我们这一对样本整体相关，即$\bm A$中至少有一个1，却不告诉我们1在哪里。对齐方法需要在训练中把$\bm A$当作未观测的潜在变量，用相似度估计它，或者干脆只在整体层面训练、不估计$\bm A$。

弱配对还可能是错的。网页替代文字常写“产品图”“点击放大”，与画面内容无关；教学视频中的旁白可能先讲下一步要做什么，再去做。Miech等人估计HowTo100M中约一半的片段与其旁白在内容上并不对齐\scite{vision-milnce2020}这类噪声有两种形态。一种是假正例：标为配对的两项其实无关。另一种是假负例：一批数据中，被当作非配对的另一条记录恰好也符合描述，例如两张图都拍到了红色汽车。前者让模型把无关内容拉近，后者让模型把相关内容推远。HowTo100M数据集本身及其检索用法见\ref{sec:vision-understanding-howto}节。

噪声的比例决定了能从数据中学到多少。用一个教学设定说明：若训练对中有比例$\epsilon$为假正例，且假正例的两项相互独立，那么按配对标签计算的平均相似度，等于真实配对平均相似度与独立配对平均相似度的加权和$(1-\epsilon)\bar s_{+}+\epsilon\bar s_{0}$。$\epsilon=0.5$时，配对信号被稀释一半，但只要$\bar s_{+}>\bar s_{0}$，增大真实配对相似度仍能降低目标，学习方向不变；需要的数据量随之增加。若假正例不是独立的，而是系统地偏向某类内容，例如旁白总在描述画面之外的工具，学到的对应就会系统地偏离。对比学习怎样在目标中容纳多个候选正例，见\ref{sec:vision-alignment-contrastive}节。

\subsection{模态的缺失与冲突}
\label{sec:vision-alignment-missing-conflict}

前两小节假定每个模态都存在且如实记录同一事件。实际数据常常违反这一点。开篇录像中，录音可能因设备故障整段缺失，也可能在第5至第7秒被风声覆盖；字幕可能根本没有生成。我们用一个掩码向量$\bm r\in\{0,1\}^{M}$记录$M$个模态中哪些可用，$r_m=1$表示第$m$个模态被观测到。缺失可以是整段的（$r_m=0$），也可以是局部的：一段音频中只有某些时间单元可用，需要按时间单元记录掩码。

缺失本身可能携带信息。若录音设备只在检测到大声时才开始录音，那么“录音缺失”就暗示周围比较安静。卷二统计推断章区分了三种缺失机制：缺失与所有数据无关时为\term{完全随机缺失}（missing completely at random, MCAR）；给定已观测变量后缺失与缺失值本身无关时为\term{随机缺失}（missing at random, MAR）；缺失依赖缺失值本身时为\term{非随机缺失}（missing not at random, MNAR）。只有前两种情形下，用已观测部分估计缺失部分才不会引入系统偏差，可参考\BookChapterRef{chap:statistical-inference}{卷2《数学准备：概率、因果与决策》的“统计推断”章}。多模态数据中的缺失常属于第三种，例如低质量视频更可能没有字幕，这使“只用完整样本训练”也会改变数据分布。

冲突是另一种情形：各模态都在，却给出不同的结论。开篇录像中，鸣笛声可能来自画面外的另一辆车；字幕说“右转”，画面里的红车却在左转，可能是语音识别错误，也可能是旁白者口误。冲突时简单平均两个模态的判断，会得到一个两边都不支持的中间结论。处理冲突需要先判断哪一侧更可信，或者显式报告不一致，而不是默认各模态说的都是同一件事。\ref{sec:vision-alignment-late-fusion}节会用一个时刻估计的算例说明，相乘式的融合在冲突时会给出过于自信的错误结果。

这三个特点分别对应后面三节。异构与弱配对使跨模态对应不能直接读出，需要\ref{sec:vision-alignment-alignment}节的方法建立；对应建立后仍需决定在哪里、用什么形式合并证据，这是\ref{sec:vision-alignment-fusion}节的内容；缺失与冲突则要求训练和推理不依赖“所有模态都在”的假设，这是\ref{sec:vision-alignment-missing}节的内容。

\section{跨模态对齐}
\label{sec:vision-alignment-alignment}

\term{跨模态对齐}（cross-modal alignment）要找出不同模态单元之间的对应关系。输入是同一事件的两种记录，例如一段视频的帧序列与一句字幕的词序列；输出可以是一个对应矩阵、一条对应路径，或者一个使对应内容彼此接近的共同表示空间。难点来自上一节：两侧单元不能直接比较，训练时多数只知道整体配对，对应关系可能是一对多、多对一，也可能有一侧单元没有对应。

对齐首先要规定粒度，也就是“对应”发生在多大的单元之间。表\ref{tab:vision-alignment-granularity}列出本卷会遇到的四种粒度。粒度越细，能回答的问题越具体，例如“鸣笛发生在哪几帧”，所需的监督或结构假设也越强。

\begin{table}[htbp]
  \centering
  \small
  \begin{tabularx}{\linewidth}{@{}p{16mm}p{30mm}XX@{}}
    \toprule
    粒度 & 对应的单元 & 常见监督 & 本卷位置\\
    \midrule
    全局 & 整图与整句；整段视频与整段音频 & 整体配对 & CLIP、ImageBind（\ref{chap:vision-multimodal-models}章）\\
    区域—短语 & 检测框或图块与名词短语 & 整体配对，或框—短语标注 & 文字短语的区域定位（\ref{chap:vision-multimodal-understanding}章）\\
    帧—词 & 视频片段与旁白中的词句 & 带时间戳的旁白，常有错位 & 跨模态检索（\ref{chap:vision-multimodal-understanding}章）\\
    音画同步 & 音频帧与视频帧 & 同一录像中的时间共现 & 音画时间定位与声音生成（\ref{chap:vision-multimodal-understanding}、\ref{chap:vision-multimodal-generation}章）\\
    \bottomrule
  \end{tabularx}
  \caption{跨模态对齐的四种粒度。同一系统可以同时使用多种粒度，例如先用全局对齐召回候选，再用区域—短语对齐检查细节。}
  \label{tab:vision-alignment-granularity}
\end{table}

对齐方法可以按它们回答的问题分成两类。一类回答“怎样让两种表示可以比较”：给定配对样本，学习两个映射，使配对的两项在映射后接近。典型相关分析用线性投影最大化相关系数，对比学习用神经网络编码器和非配对样本作对照。另一类回答“可比较之后，哪个单元对应哪个单元”：给定两条序列和单元间的相似度，求一个对应关系。动态时间规整要求对应保持时间顺序，用动态规划求代价最小的硬对应；注意力不强制顺序，用归一化权重给出可微的软对应。实际系统常把两类组合：先用对比学习得到可比较的特征，再用动态时间规整或注意力在特征上求对应。下面四个小节依次讨论这四种方法，最后比较它们的假设与代价。

\subsection{基于典型相关分析的对齐}
\label{sec:vision-alignment-cca}

最简单的情形是：每个样本有两组已经提取好的特征向量，例如视觉特征$\bm x\in\mathbb R^{d_x}$与音频特征$\bm y\in\mathbb R^{d_y}$，我们希望各找一个方向，使两组特征在这两个方向上的投影尽可能同步变化。若红车出现时视觉特征的某个组合变大、发动机声的某个组合也变大，这两个组合就是我们要找的共同成分。分别对两侧做主成分分析不能完成这件事：主成分分析选的是各自方差最大的方向，视觉一侧方差最大的变化可能是光照，与声音无关。

\term{典型相关分析}（canonical correlation analysis, CCA）改为直接最大化两侧投影之间的相关系数\scite{hotelling1936cca}设特征已经中心化，两侧协方差为$\bm\Sigma_{xx}=\mathbb E[\bm x\bm x^\top]$、$\bm\Sigma_{yy}=\mathbb E[\bm y\bm y^\top]$，交叉协方差为$\bm\Sigma_{xy}=\mathbb E[\bm x\bm y^\top]=\bm\Sigma_{yx}^\top$。投影$u=\bm w_x^\top\bm x$与$v=\bm w_y^\top\bm y$的相关系数为
\begin{equation}
  \rho(\bm w_x,\bm w_y)=
  \frac{\bm w_x^\top\bm\Sigma_{xy}\bm w_y}
       {\sqrt{\bm w_x^\top\bm\Sigma_{xx}\bm w_x}\,
        \sqrt{\bm w_y^\top\bm\Sigma_{yy}\bm w_y}}.
  \label{eq:vision-alignment-cca-objective}
\end{equation}
卷四降维章给出了这个问题的一种解法：先分别白化两侧，再对白化后的交叉协方差$\bm\Sigma_{xx}^{-1/2}\bm\Sigma_{xy}\bm\Sigma_{yy}^{-1/2}$作奇异值分解，奇异值就是典型相关系数，可参考\BookChapterRef{chap:dimensionality-reduction}{卷四“降维”章}。这里换一条路，用拉格朗日乘子把它推成广义特征值问题，因为这种形式更直接地显示正则化和非线性推广改动了哪一项。

缩放$\bm w_x$或$\bm w_y$不改变式\eqref{eq:vision-alignment-cca-objective}，所以可以固定两侧投影的方差为1，把问题改写为
\[
  \max_{\bm w_x,\bm w_y}\ \bm w_x^\top\bm\Sigma_{xy}\bm w_y
  \quad\text{满足}\quad
  \bm w_x^\top\bm\Sigma_{xx}\bm w_x=1,\ \
  \bm w_y^\top\bm\Sigma_{yy}\bm w_y=1.
\]
这一改写要求$\bm\Sigma_{xx}$、$\bm\Sigma_{yy}$正定，否则存在方差为零的方向，相关系数无定义。引入乘子$\lambda_x,\lambda_y$，拉格朗日函数为
\[
  \mathcal L=\bm w_x^\top\bm\Sigma_{xy}\bm w_y
   -\frac{\lambda_x}{2}\bigl(\bm w_x^\top\bm\Sigma_{xx}\bm w_x-1\bigr)
   -\frac{\lambda_y}{2}\bigl(\bm w_y^\top\bm\Sigma_{yy}\bm w_y-1\bigr).
\]
对$\bm w_x$与$\bm w_y$求梯度并令其为零，利用$\bm\Sigma_{xx}$、$\bm\Sigma_{yy}$对称，得到
\begin{equation}
  \bm\Sigma_{xy}\bm w_y=\lambda_x\bm\Sigma_{xx}\bm w_x,\qquad
  \bm\Sigma_{yx}\bm w_x=\lambda_y\bm\Sigma_{yy}\bm w_y.
  \label{eq:vision-alignment-cca-stationary}
\end{equation}
用$\bm w_x^\top$左乘第一式，并代入约束$\bm w_x^\top\bm\Sigma_{xx}\bm w_x=1$，得到$\lambda_x=\bm w_x^\top\bm\Sigma_{xy}\bm w_y$；同理用$\bm w_y^\top$左乘第二式得到$\lambda_y=\bm w_y^\top\bm\Sigma_{yx}\bm w_x$。两者是同一个标量的转置，故$\lambda_x=\lambda_y$，而且它就是驻点处的目标值，也就是相关系数$\rho$。把两式合写成分块形式：
\begin{equation}
  \begin{bmatrix}\bm 0&\bm\Sigma_{xy}\\ \bm\Sigma_{yx}&\bm 0\end{bmatrix}
  \begin{bmatrix}\bm w_x\\ \bm w_y\end{bmatrix}
  =\rho
  \begin{bmatrix}\bm\Sigma_{xx}&\bm 0\\ \bm 0&\bm\Sigma_{yy}\end{bmatrix}
  \begin{bmatrix}\bm w_x\\ \bm w_y\end{bmatrix}.
  \label{eq:vision-alignment-cca-gep}
\end{equation}
这是对称矩阵与对称正定矩阵构成的\term{广义特征值问题}（generalized eigenvalue problem）。卷一线性空间章说明了这类问题的性质：右侧正定时，广义特征值全为实数，通过Cholesky换元可化为普通对称特征值问题，最大的广义特征值就是广义Rayleigh商$\bm z^\top\bm A\bm z/\bm z^\top\bm B\bm z$的最大值，可参考\BookChapterRef{chap:linear-analysis}{卷一“特征值与谱理论”章}。因此式\eqref{eq:vision-alignment-cca-gep}的最大广义特征值就是第一典型相关系数，对应的广义特征向量给出两侧方向。

维数较小时，也可以消去一侧。$\rho\ne0$时，由第二式得$\bm w_y=\bm\Sigma_{yy}^{-1}\bm\Sigma_{yx}\bm w_x/\rho$，代入第一式：
\begin{equation}
  \bm\Sigma_{xy}\bm\Sigma_{yy}^{-1}\bm\Sigma_{yx}\bm w_x
  =\rho^2\bm\Sigma_{xx}\bm w_x.
  \label{eq:vision-alignment-cca-reduced}
\end{equation}
这是一个$d_x$维的广义特征值问题，特征值是相关系数的平方。求得$\bm w_x$后再按上式算出$\bm w_y$。第二对、第三对方向取后续的广义特征向量；它们满足$\bm w_{x,k}^\top\bm\Sigma_{xx}\bm w_{x,l}=0$（$k\ne l$），即各对投影之间互不相关，这是协方差度量下的正交，不一定是欧氏正交。

\begin{example}[两维特征上的典型相关]
  \label{ex:vision-alignment-cca}
  以下为教学构造。设视觉特征$\bm x=(x_1,x_2)^\top$、音频特征$\bm y=(y_1,y_2)^\top$的协方差为
  \[
    \bm\Sigma_{xx}=\begin{bmatrix}4&0\\0&1\end{bmatrix},\quad
    \bm\Sigma_{yy}=\begin{bmatrix}1&0\\0&1\end{bmatrix},\quad
    \bm\Sigma_{xy}=\begin{bmatrix}0.4&0\\0&0.8\end{bmatrix}.
  \]
  $x_1$可看作随光照变化的亮度成分，方差最大；$x_2$可看作车辆运动强度，与音频能量$y_2$同步。对视觉一侧做主成分分析，会选方差为4的$x_1$。按式\eqref{eq:vision-alignment-cca-reduced}，左侧矩阵为$\bm\Sigma_{xy}\bm\Sigma_{yx}=\operatorname{diag}(0.16,0.64)$，与$\bm\Sigma_{xx}$都是对角阵，广义特征值逐坐标相除得$0.16/4=0.04$与$0.64/1=0.64$，相关系数为$0.2$与$0.8$。第一对方向因此是$\bm w_x=(0,1)^\top$，满足$\bm w_x^\top\bm\Sigma_{xx}\bm w_x=1$；代回得$\bm w_y=\bm\Sigma_{yx}\bm w_x/0.8=(0,1)^\top$。CCA选出了方差较小却与声音同步的运动成分，主成分分析选出的亮度成分与音频的相关系数只有$0.4/(2\times1)=0.2$。
\end{example}

用样本估计协方差时，CCA容易把偶然的相关当成共同成分。设有$n$个中心化样本，把视觉特征排成矩阵$\bm X\in\mathbb R^{n\times d_x}$，音频特征排成$\bm Y\in\mathbb R^{n\times d_y}$。投影得分$\bm X\bm w_x$落在$\bm X$的列空间中；中心化使每列元素之和为零，所以列空间最多是$n-1$维。若$d_x\geq n-1$且样本处于一般位置，$\bm X$的列空间就是全部和为零的$n$维向量，任意$\bm Y\bm w_y$都能被某个$\bm X\bm w_x$精确表示，样本相关系数等于1。例如$n=3$，三个视觉样本为$(1,0)$、$(0,1)$、$(-1,-1)$，音频得分为$(2,-1,-1)$，取$\bm w_x=(2,-1)^\top$就得到完全相同的得分向量。这个相关系数1不说明两侧有任何共同成分，只说明样本太少。

正则化的做法是把式\eqref{eq:vision-alignment-cca-gep}右侧换成$\bm\Sigma_{xx}+r_x\bm I$与$\bm\Sigma_{yy}+r_y\bm I$，$r_x,r_y>0$。这保证右侧正定，也惩罚只靠小方差方向凑出相关的投影。仍用例\ref{ex:vision-alignment-cca}，取$r_x=r_y=0.5$，式\eqref{eq:vision-alignment-cca-reduced}变为$\operatorname{diag}(0.16,0.64)/1.5$与$\operatorname{diag}(4.5,1.5)$的广义特征值问题，两个值约为$0.0237$与$0.2844$，开方约为$0.154$与$0.533$。方向未变，但估计的相关系数都被压低；这时它们不再是原来意义上的相关系数，只用于排序方向。正则化强度通常在留出的配对样本上选择。

CCA只能发现线性的共同成分。若声音能量与运动速度的平方成正比，线性投影的相关系数可能很低。\term{核典型相关分析}（kernel canonical correlation analysis）先用核函数把两侧映射到高维特征空间再求相关方向；\term{深度典型相关分析}（deep canonical correlation analysis）把两侧换成神经网络$f_{\bm\theta_x}(\bm x)$与$g_{\bm\theta_y}(\bm y)$，以网络输出上的总典型相关为目标，反向传播联合训练两个网络\scite{andrew2013dcca}两种推广在卷四降维章中都有说明。深度CCA与后面的对比学习目标相近，区别在于防止退化的方式：CCA用方差约束$\bm w^\top\bm\Sigma\bm w=1$阻止两侧把所有样本映射到同一点，对比学习则用非配对样本作对照。

CCA作为对齐方法有三个限制。它只建立全局的共同空间，不说明一段视频的哪一帧对应哪个词；它要求每个样本两侧都有固定维数的特征，变长序列需要先汇总；它最大化的是相关性，不考虑哪些成分对下游任务有用，单侧独有却很重要的信息会被丢掉。这些限制分别由后面三个小节在不同方向上处理。

\subsection{基于动态时间规整的对齐}
\label{sec:vision-alignment-dtw}

CCA给出的是整段样本之间的对应。开篇问题还需要更细的对应：录音中的鸣响对应哪几帧画面。若两条序列以相同速度记录同一过程，只需估计一个固定的时间偏移，\ref{sec:vision-understanding-av-time}节的滑动窗口比较就是这样做的。困难在于速度往往不同：同一个转弯动作，两次录像中可能一次用2秒、一次用4秒；讲解视频中的旁白语速与操作节奏也不一致。把一条序列均匀拉伸到另一条的长度，会在快慢变化处错位。

\term{动态时间规整}（dynamic time warping, DTW）允许两条序列以不同的局部速度对应，同时保持时间先后顺序\scite{vision-sakoe1978}设两条特征序列为$\bm x_1,\ldots,\bm x_n$和$\bm y_1,\ldots,\bm y_m$，$c(i,j)=d(\bm x_i,\bm y_j)$是第$i$个与第$j$个单元之间的局部代价，例如欧氏距离。一条\term{规整路径}（warping path）是格点序列$P=((i_1,j_1),\ldots,(i_L,j_L))$，它满足三个条件：边界条件$(i_1,j_1)=(1,1)$、$(i_L,j_L)=(n,m)$，要求两条序列的首尾互相对应；单调条件$i_{l+1}\geq i_l$、$j_{l+1}\geq j_l$，要求对应不能倒退；连续条件$(i_{l+1}-i_l,\,j_{l+1}-j_l)\in\{(1,0),(0,1),(1,1)\}$，要求每步只前进一格，不能跳过单元。路径长度$L$在$\max(n,m)$与$n+m-1$之间。DTW代价是所有规整路径中局部代价之和的最小值：
\begin{equation}
  \operatorname{DTW}(\bm x,\bm y)=\min_{P}\sum_{l=1}^{L}c(i_l,j_l).
  \label{eq:vision-alignment-dtw}
\end{equation}
水平或竖直的一步表示一侧停在原单元、另一侧前进，也就是一个单元对应多个单元；斜向的一步表示两侧同步前进。

规整路径的数目随$n,m$指数增长，不能逐条枚举。把路径看成网格图上从$(1,1)$到$(n,m)$的有向路，每个格点的权重是$c(i,j)$，DTW就是这张无环图上的最短路问题。卷一图论章说明，无环图上的最短路满足最优子结构：最短路的任一前缀也是到达该前缀终点的最短路，因此可以按拓扑顺序递推，可参考\BookChapterRef{chap:graph-theory}{卷一“图的表示与计算”章}。这里的论证可以直接写出。记$D(i,j)$为从$(1,1)$到$(i,j)$的最小累计代价。由连续条件，任何到达$(i,j)$的路径，前一格必为$(i-1,j)$、$(i,j-1)$、$(i-1,j-1)$之一；若路径在前一格之前的部分不是到达前一格的最优路径，换成最优路径后总代价会降低，与最优矛盾。于是
\begin{equation}
  D(i,j)=c(i,j)+\min\bigl\{D(i-1,j),\,D(i,j-1),\,D(i-1,j-1)\bigr\},
  \label{eq:vision-alignment-dtw-recursion}
\end{equation}
边界为$D(1,1)=c(1,1)$，并约定$D(0,j)=D(i,0)=+\infty$，使第一行、第一列只能从一个方向到达。$D(n,m)$就是式\eqref{eq:vision-alignment-dtw}的值。从$(n,m)$出发，每次回到取得最小值的前一格，直到$(1,1)$，就得到最优路径。算法\ref{alg:vision-alignment-dtw}写出完整步骤。

\begin{algorithm}[动态时间规整]
  \label{alg:vision-alignment-dtw}
  \begin{algorithmic}[1]
    \Require 序列$\bm x_1,\ldots,\bm x_n$与$\bm y_1,\ldots,\bm y_m$，局部代价$d$，可选带宽$r$
    \Ensure DTW代价$D(n,m)$与规整路径$P$
    \State 令$D(i,j)\gets+\infty$，$0\leq i\leq n$，$0\leq j\leq m$；令$D(0,0)\gets0$
    \For{$i=1,\ldots,n$}
      \For{$j=1,\ldots,m$}
        \If{给定带宽且$\lvert i-jn/m\rvert>r$}
          \State 跳过此格，保持$D(i,j)=+\infty$
        \Else
          \State $D(i,j)\gets d(\bm x_i,\bm y_j)+\min\{D(i-1,j),D(i,j-1),D(i-1,j-1)\}$
        \EndIf
      \EndFor
    \EndFor
    \State 令$(i,j)\gets(n,m)$，$P\gets((n,m))$
    \While{$(i,j)\ne(1,1)$}
      \State 在$(i-1,j)$、$(i,j-1)$、$(i-1,j-1)$中选$D$最小者为新的$(i,j)$，并插入$P$前端
    \EndWhile
    \State \Return $D(n,m)$与$P$
  \end{algorithmic}
\end{algorithm}

初始化$D(0,0)=0$与$D(1,1)=c(1,1)$等价，因为$(1,1)$的三个前驱中只有$(0,0)$有限。双重循环逐格计算一次，时间与存储都是$O(nm)$；只要代价值而不要路径时，按行递推只需保存两行，存储降为$O(m)$。

\begin{example}[动作强度与音频能量的规整]
  \label{ex:vision-alignment-dtw}
  以下为教学构造。设开篇录像中5帧画面的动作强度为$x=(0,0,2,4,4)$，3个音频单元的能量为$y=(0,2,4)$，局部代价$c(i,j)=\lvert x_i-y_j\rvert$。按式\eqref{eq:vision-alignment-dtw-recursion}逐行计算。第一列只能从上方到达：$D(1,1)=0$，$D(2,1)=0+0=0$，$D(3,1)=2+0=2$，$D(4,1)=4+2=6$，$D(5,1)=4+6=10$。第二列：$D(1,2)=2+0=2$；$D(2,2)=2+\min\{2,0,0\}=2$；$D(3,2)=0+\min\{2,2,0\}=0$；$D(4,2)=2+\min\{0,6,2\}=2$；$D(5,2)=2+\min\{2,10,6\}=4$。第三列：$D(1,3)=4+2=6$；$D(2,3)=4+\min\{6,2,2\}=6$；$D(3,3)=2+\min\{6,0,2\}=2$；$D(4,3)=0+\min\{2,2,0\}=0$；$D(5,3)=0+\min\{0,4,2\}=0$。回溯得到路径$(1,1),(2,1),(3,2),(4,3),(5,3)$，总代价为0：前两帧都对应静止的第一个音频单元，最后两帧都对应最响的第三个单元。若把5帧均匀映射到3个单元，第$i$帧对应第$\lceil 3i/5\rceil$个单元，即$(1,1),(2,2),(3,2),(4,3),(5,3)$，第2帧的静止画面被配给能量为2的单元，总代价为2。
\end{example}

图\ref{fig:vision-alignment-dtw-path}画出了这张累计代价表与回溯路径。读图时沿箭头看：水平箭头表示同一个音频单元持续对应多帧画面，斜向箭头表示两侧同步前进。路径上每格的累计代价都是0，说明每一步都找到了完全匹配的单元；均匀映射没有这种灵活性。

\begin{figure}[htbp]
  \centering
  % DTW教学算例的累计代价表；宽约85mm。x=(0,0,2,4,4)，y=(0,2,4)，局部代价|x_i-y_j|。
\begin{tikzpicture}[x=1.25cm,y=1.0cm,font=\figurefont\small,
  arr/.style={-{Stealth[length=2mm]},BookTeal,line width=1.0pt}]
  % 网格：列为视频帧i=1..5，行为音频单元j=1..3
  \foreach \i/\xv in {1/0,2/0,3/2,4/4,5/4}{
    \node at (\i,-0.05) {$x_{\i}=\xv$};
  }
  \foreach \j/\yv in {1/0,2/2,3/4}{
    \node[anchor=east] at (0.4,\j) {$y_{\j}=\yv$};
  }
  \foreach \i/\j in {1/1,2/1,3/2,4/3,5/3}{
    \fill[BookTeal!15] (\i-0.5,\j-0.45) rectangle (\i+0.5,\j+0.45);
  }
  \foreach \i in {1,...,5}{
    \foreach \j in {1,2,3}{
      \draw[BookRule,line width=0.5pt] (\i-0.5,\j-0.45) rectangle (\i+0.5,\j+0.45);
    }
  }
  \foreach \i/\j/\v in {1/1/0,1/2/2,1/3/6,
                        2/1/0,2/2/2,2/3/6,
                        3/1/2,3/2/0,3/3/2,
                        4/1/6,4/2/2,4/3/0,
                        5/1/10,5/2/4,5/3/0}{
    \node at (\i,\j) {$\v$};
  }
  % 回溯得到的规整路径
  \draw[arr] (1.2,1) -- (1.8,1);
  \draw[arr] (2.2,1.15) -- (2.8,1.85);
  \draw[arr] (3.2,2.15) -- (3.8,2.85);
  \draw[arr] (4.2,3) -- (4.8,3);
  \node[anchor=north] at (3,-0.4) {视频帧动作强度$x_i$};
  \node[rotate=90,anchor=south] at (-0.75,2) {音频能量$y_j$};
\end{tikzpicture}

  \caption{例\ref{ex:vision-alignment-dtw}的累计代价表。横轴是5帧画面的动作强度，纵轴是3个音频单元的能量，格中数字为$D(i,j)$；浅色格与箭头是从$(5,3)$回溯得到的规整路径。数值为教学构造，局部代价为绝对差。}
  \label{fig:vision-alignment-dtw-path}
\end{figure}

不加约束的DTW可能给出不合理的对应。若某一单元的代价特别低，路径可能在那里长时间停留，让一侧的一个单元对应另一侧的一大段，这称为病态规整。Sakoe与Chiba为此提出两类约束：限制路径只能在对角线附近的带状区域内，即$\lvert i-j\rvert\leq r$；以及限制路径连续水平或竖直移动的步数，即限制斜率。两条序列长度不同时，带状区域按对角线比例改写为$\lvert i-jn/m\rvert\leq r$，算法\ref{alg:vision-alignment-dtw}中的带宽即此意。带宽同时把计算量降到$O(\max(n,m)\,r)$。代价是，若真实的速度差超出带宽，最优路径就被排除在外。

DTW作为对齐方法有几个需要注意的条件。它要求对应关系单调，若旁白先讲第三步再讲第一步，单调路径无法表示这种倒序。它要求首尾对应，若一段视频只有中间部分与音频对应，需要改用允许任意起止点的子序列版本。DTW代价随路径长度累加，比较不同长度的序列对时，常按路径长度或$n+m$归一化。DTW也不满足三角不等式，不是严格意义上的距离。最重要的是，DTW只负责“在给定局部代价下找对应”，局部代价$d(\bm x_i,\bm y_j)$是否反映语义对应，取决于两侧特征是否已经对齐到可比较的空间；跨模态使用时，常先用CCA或对比学习得到特征，再在其上运行DTW。DTW的离散路径也不可微，难以作为网络中间层端到端训练，这引出下一小节的软对齐。

\subsection{基于注意力的软对齐}
\label{sec:vision-alignment-attention}

DTW的输出是一条硬路径：每一对单元要么对应，要么不对应。在神经网络中，我们希望对齐关系能随网络参数一起学习，而硬路径的选择对参数不可微，误差无法通过“选哪一格”传回特征提取部分。此外，单调约束并非总是成立：一句描述中的词序与图像中区域的位置之间没有先后关系，“红色汽车右侧的自行车”先提汽车后提自行车，图中自行车却可能在左边。

注意力机制给出另一种对应方式。对目标序列的每个单元，它计算与源序列每个单元的相关分数，再用softmax归一化成权重，按权重读取源单元的内容。卷五Transformer章说明了这一运算，可参考\BookChapterRef{chap:transformer}{卷五“Transformer”章}。这里关心的是它的权重矩阵：设目标一侧有$m$个单元，投影后的查询为$\bm Q\in\mathbb R^{m\times d_{\mathrm k}}$，源一侧有$n$个单元，键为$\bm K\in\mathbb R^{n\times d_{\mathrm k}}$、值为$\bm V\in\mathbb R^{n\times d_{\mathrm o}}$，则
\begin{equation}
  \bm A=\operatorname{softmax}_{\mathrm{row}}\!\left(\frac{\bm Q\bm K^\top}{\sqrt{d_{\mathrm k}}}\right)\in\mathbb R^{m\times n},
  \qquad
  \bm R=\bm A\bm V.
  \label{eq:vision-alignment-attention}
\end{equation}
$\bm A$的每行非负且和为1，$A_{i,j}$可以读作“目标第$i$个单元在多大程度上对应源第$j$个单元”。与\ref{sec:vision-alignment-weak-pairing}节的0/1对齐矩阵相比，$\bm A$是一个\term{软对齐}（soft alignment）：它把一个目标单元的对应分摊到多个源单元上，而且整个运算可微，对齐关系可以由下游目标间接训练，不需要对齐标注。Bahdanau等人在机器翻译中引入注意力时，正是把它解释为“联合学习对齐与翻译”\scite{bahdanau2015}这种跨序列的读取称为\term{交叉注意力}（cross-attention），视觉语言模型中读取视觉特征的写法见式\eqref{eq:vision-models-cross-read}。

用一个小例子看软对齐矩阵怎样形成。设字幕有两个词“红色汽车”“鸣笛”，图像有3个区域：右侧红车、左侧自行车、路面。假设缩放后的相关分数为
\[
  \frac{\bm Q\bm K^\top}{\sqrt{d_{\mathrm k}}}=
  \begin{bmatrix}2&0&1\\0&1&1\end{bmatrix},
  \qquad
  \bm A\approx
  \begin{bmatrix}0.665&0.090&0.245\\0.155&0.422&0.422\end{bmatrix}.
\]
第一行把约三分之二的权重给了红车区域，可以说“红色汽车”软对齐到红车；第二行在自行车和路面之间平分，说明“鸣笛”在静态画面中找不到明确的视觉对应。这种分散正是需要关注的情形：softmax强制每行权重和为1，即使没有任何区域真正对应，权重也必须分给某些区域。注意力权重因此只说明模型内部怎样加权读取，不能直接当作定位证据。

区域—短语粒度的对齐可以由整体配对训练出来。Karpathy与Fei-Fei用一个检测器从每张图取出19个候选框，加上整图共20个区域向量$\bm v_k$，用双向循环网络把句子中每个词编码为$\bm s_t$，两者在同一维数的空间中\scite{vision-karpathy2015}图像$p$与句子$q$的整体分数定义为
\begin{equation}
  S_{p,q}=\sum_{t\in q}\max_{k\in p}\bm v_k^\top\bm s_t,
  \label{eq:vision-alignment-karpathy}
\end{equation}
即每个词只取与它最匹配的区域，再对所有词求和。训练只用整体配对：要求配对的图句分数比同批中非配对的高出间隔1，用合页损失惩罚违反的情形。由于整体分数由词—区域的最大内积构成，提高配对分数的唯一途径是让每个词找到一个匹配的区域，区域—短语对应就作为副产品形成。在上面的算例中，式\eqref{eq:vision-alignment-karpathy}的分数为$2+1=3$。与软注意力相比，这里用$\max$取了硬对应，但梯度仍可沿被选中的那一项传回。原文还在词的对应上加了链状马尔可夫随机场，用动态规划鼓励相邻的词对应同一区域，使对应落在连续的短语上。区域—短语对齐的检测式做法见\ref{sec:vision-understanding-mdetr}节的MDETR。

对时间序列，注意力的优点是不需要预先把两侧重采样到同一帧率。Tsai等人的多模态Transformer（multimodal transformer, MulT）处理文字、视觉、声学三种未对齐序列，为每一对有序模态设一个跨模态Transformer，共6个；在从模态$\beta$到模态$\alpha$的模块中，查询来自$\alpha$，键和值来自$\beta$，输出长度与$\alpha$相同\scite{vision-mult2019}该文所用数据的声学与视觉特征分别约每秒12.5帧和15帧，文字按词排列，三者长度各不相同，注意力权重自行决定哪些单元相关。

不加约束的注意力也有代价。它不知道时间先后，一个词可以同时对应开头和结尾；序列越长，权重越容易分散。语音识别、同步配音这类任务中对应关系本来就是单调的，忽略这一结构既浪费数据，也可能产生跳跃或重复的对齐。Raffel等人的单调注意力把DTW的单调约束带回到可微的框架中\scite{vision-monotonic2017}推理时，对第$i$个目标单元，模型从上一步停下的源位置开始向后扫描，在每个位置$j$以概率$p_{i,j}=\sigma(e_{i,j})$决定“停在这里”，$\sigma$是logistic函数，$e_{i,j}$是相关分数；停下的位置就是本步对应的源单元，下一步只能从这里继续向后。训练时不采样，而是计算停在每个位置的概率$\alpha_{i,j}$作为软对齐权重。停在$(i,j)$要求：本步到达位置$j$且在此停下。到达$j$有两种途径：上一步恰好停在$j$，或本步从$j-1$经过而没有停下。记$q_{i,j}$为本步到达位置$j$的概率，则
\begin{equation}
  q_{i,j}=(1-p_{i,j-1})\,q_{i,j-1}+\alpha_{i-1,j},\qquad
  \alpha_{i,j}=p_{i,j}\,q_{i,j},
  \label{eq:vision-alignment-monotonic}
\end{equation}
边界为$p_{i,0}=0$、$q_{i,0}=0$。第一项是从前一位置经过而未停的概率，第二项是上一步停在本位置的概率。

用三个源单元走一遍。设初始对应集中在第一个源单元，即$\alpha_{0,\cdot}=(1,0,0)$。第一步的停止概率为$p_{1,\cdot}=(0.5,0.8,0.9)$：$q_{1,1}=1$，$\alpha_{1,1}=0.5$；$q_{1,2}=0.5\times1=0.5$，$\alpha_{1,2}=0.4$；$q_{1,3}=0.2\times0.5=0.1$，$\alpha_{1,3}=0.09$。余下的0.01是扫描到末尾也没有停下的概率。第二步取$p_{2,\cdot}=(0.1,0.5,0.9)$：$q_{2,1}=0.5$，$\alpha_{2,1}=0.05$；$q_{2,2}=0.9\times0.5+0.4=0.85$，$\alpha_{2,2}=0.425$；$q_{2,3}=0.5\times0.85+0.09=0.515$，$\alpha_{2,3}\approx0.464$。第二步的权重整体向右移动，而第一个源单元上只剩下上一步停在那里、本步又立刻停下的那一小部分，模型不可能把大量权重重新放回已经越过的位置。推理时扫描从上一步位置继续，总共只需检查$O(n+m)$个位置，可以随输入到达在线进行；训练时计算全部$\alpha_{i,j}$仍需$O(nm)$。

三种注意力式对齐的差别在于对对应结构作了多少假设。普通交叉注意力不作假设，适合区域与词这类没有顺序的对应，代价是权重可能分散、需要大量数据才能学准。式\eqref{eq:vision-alignment-karpathy}用最大值取硬对应，让整体配对就能训练出局部对应，但每个词只能对应一个区域。单调注意力加入时间顺序，适合语音与文字、音频与视频这类保持先后的对应，代价是不能表示倒序。多模态模型中的交叉注意力连接器，例如\ref{sec:vision-models-blip2}节的Q-Former与\ref{sec:vision-models-flamingo}节的门控交叉注意力，都属于第一种；它们的权重是否真的落在相关区域上，需要\ref{sec:vision-understanding-input-protocol}节那样的替换实验来检查。

\subsection{基于对比学习的对齐}
\label{sec:vision-alignment-contrastive}

前三小节都假设已经有了两侧可以比较的特征，或者像CCA那样只能学习线性投影。要从像素和词元出发、用数亿个弱配对学习两个深层编码器，需要一个能在大规模数据上稳定训练、又不会让所有样本塌缩到同一点的目标。\term{对比学习}（contrastive learning）用非配对样本作对照来提供这一约束：配对的两项要比同批其他组合更相似。卷八基于数据本身的知识获取章从自监督角度介绍了对比目标与负样本采样，可参考\BookChapterRef{chap:data-derived-knowledge-acquisition}{卷八“基于数据本身的知识获取”章}。这里从对齐的角度说明这一目标在优化什么、温度和负样本数怎样影响结果。

设图像编码器$f$与文字编码器$g$都输出单位向量，相似度为余弦$s(\bm x,q)=f(\bm x)^\top g(q)$。给定一个锚点$\bm x$及$K$个候选文本$q_1,\ldots,q_K$，其中$q_1$与$\bm x$配对，其余$K-1$个来自其他样本。\term{信息噪声对比估计}（information noise-contrastive estimation, InfoNCE）损失把“从$K$个候选中认出配对项”当作$K$类分类\scite{vision-cpc2018}
\begin{equation}
  L_K=-\mathbb E\left[\log
   \frac{\exp\bigl(s(\bm x,q_1)/\tau\bigr)}
        {\sum_{j=1}^{K}\exp\bigl(s(\bm x,q_j)/\tau\bigr)}\right].
  \label{eq:vision-alignment-infonce}
\end{equation}
温度$\tau>0$把相似度放大成分类分数。\ref{sec:vision-models-contrastive}节的CLIP对一批$n$对图文同时以每幅图和每段文字为锚点，把行、列两个方向的这一损失取平均，此时每个锚点的候选数$K=n$。

这个损失与\term{互信息}（mutual information）有确定的关系。卷二概率分布章证明：若配对项按联合分布抽取、其余候选彼此独立并按边缘分布抽取，对任意评分函数都有
\begin{equation}
  \log K-L_K\leq I(\bm x;q),
  \label{eq:vision-alignment-infonce-bound}
\end{equation}
而且使$L_K$最小的评分函数是对数密度比$\log[p(\bm x,q)/(p(\bm x)p(q))]$加上只依赖$\bm x$的常数，可参考\BookChapterRef{chap:information-theory-and-statistical-geometry}{卷2《数学准备：概率、因果与决策》的“概率分布的分析与优化”章}。这给出两个对齐上的解读。第一，最优评分衡量的是“这两项一同出现的可能性比随机凑在一起高多少”，这正是跨模态对应在统计上的含义；对比学习训练出的相似度，在理想情况下逼近这个比值。第二，式\eqref{eq:vision-alignment-infonce-bound}左侧不超过$\log K$，因为$L_K\geq0$。候选数为$K=3$时，无论两种模态共享多少信息，这个估计都不超过$\log3\approx1.10$奈特；CLIP的批大小32768对应约$10.40$奈特。批越大，下界越可能反映真实的互信息，这是对比学习偏好大批次的一个理由；但下界变紧不等于表示变好，互信息相同的表示可以对下游任务有用程度不同。

温度控制梯度在负样本之间怎样分配。记$p_j$为式\eqref{eq:vision-alignment-infonce}中第$j$个候选的softmax概率，对单个锚点的损失$\ell=-\log p_1$求导，与\ref{sec:vision-models-contrastive}节的推导相同，得到
\begin{equation}
  \frac{\partial\ell}{\partial s(\bm x,q_j)}=\frac{p_j-\mathbb I[j=1]}{\tau}.
  \label{eq:vision-alignment-infonce-grad}
\end{equation}
负样本$j$受到的推力与$p_j$成正比，两个负样本推力之比为$p_j/p_k=\exp\bigl((s_j-s_k)/\tau\bigr)$。$\tau$越小，这个比值对相似度差越敏感，梯度越集中在与锚点最相似的负样本上，即\term{难负例}（hard negative）。

\begin{example}[温度对负样本梯度的分配]
  \label{ex:vision-alignment-temperature}
  以下为教学构造。锚点图像与三个候选文本的余弦相似度为$(0.8,0.6,0.1)$，第一个是配对文本，第二个是描述相似场景的难负例，第三个是无关文本。取$\tau=1$时，softmax概率约为$(0.432,0.354,0.215)$，损失约$0.840$，式\eqref{eq:vision-alignment-infonce-grad}给出的梯度约为$(-0.568,0.354,0.215)$，难负例与无关负例受到的推力之比约为$1.65$。取$\tau=0.1$时，概率约为$(0.880,0.119,0.0008)$，损失约$0.128$，梯度约为$(-1.199,1.191,0.008)$，推力之比约为$e^{5}\approx148$，几乎全部落在难负例上。在这一组候选上，式\eqref{eq:vision-alignment-infonce-bound}左侧的单组估计分别约为$\log3-0.840\approx0.259$与$\log3-0.128\approx0.971$奈特。单组数值不是期望，不能据此比较两个温度下的互信息估计。
\end{example}

小温度让模型专注于区分细微差别，这在负样本确实不相关时有利。弱配对数据中的难负例却可能是假负例：同批中另一张图也拍到了红色汽车，它的描述与锚点几乎同义。小温度会把这类样本当作最该推开的对象，加剧\ref{sec:vision-alignment-weak-pairing}节所说的假负例伤害。CLIP把温度设为可学习参数，由训练数据决定这一折中。增大批次也有类似的两面性：负样本越多，越可能包含难负例，也越可能包含假负例。

对假正例，可以放宽“一个锚点只有一个配对”的假设。MIL-NCE针对教学视频中旁白与画面错位的问题，把一个视频片段在时间上最近的若干段旁白都当作候选正例，组成集合$\mathcal P$，同批其他样本的旁白组成负例集合$\mathcal N$\scite{vision-milnce2020}损失为
\begin{equation}
  L=-\log
  \frac{\sum_{q\in\mathcal P}\exp\bigl(s(\bm x,q)\bigr)}
       {\sum_{q\in\mathcal P}\exp\bigl(s(\bm x,q)\bigr)
        +\sum_{q'\in\mathcal N}\exp\bigl(s(\bm x,q')\bigr)}.
  \label{eq:vision-alignment-milnce}
\end{equation}
分子对候选正例求和，只要其中有一个与画面真正对应，就能把分子抬高；无关的候选正例受到的拉力较小：对某个候选正例分数的梯度与它在候选正例内部的softmax权重成正比，梯度主要流向本来就较相似的那一个。这是\term{多示例学习}（multiple-instance learning）的思路：只知道一组示例中至少有一个是正的。原文主实验使用5个候选正例，候选过多会纳入时间上较远的无关旁白。

对比学习可以在表\ref{tab:vision-alignment-granularity}的不同粒度上构造正负例。全局粒度以整图与整句为配对，CLIP与ImageBind都是这样（\ref{sec:vision-models-imagebind}节）。区域—短语粒度可以把整体分数定义为局部匹配的汇总，例如式\eqref{eq:vision-alignment-karpathy}的逐词最大值，或BLIP-2中取32个查询输出与文本相似度的最大值，再把这个整体分数放进对比损失。帧—词粒度如MIL-NCE，以时间邻近的旁白为正例。音画同步粒度可以把同一视频中时间正确对齐的音频片段当作正例，把时间错开的片段当作负例；这样的负例与正例画面相同、只差时间，迫使模型学到时间上的对应，而不只是“这类画面配这类声音”。\ref{sec:vision-multigen-audio-alignment}节中的同步特征即服务于这一需要。

对比学习与CCA要解决的是同一个问题，都是让两侧映射后的配对项接近，区别在于怎样阻止平凡解。若只要求配对项接近，把所有输入映射到同一个向量就能达到目标。CCA用方差归一化约束排除这种解，对比学习用负样本：所有输入映射到同一点时，配对项与负样本无法区分，$L_K=\log K$达到最大。对比学习因此能训练任意深的编码器，代价是需要大量负样本和与之相伴的计算，以及前面讨论的假负例问题。它得到的仍是一个可比较的空间，不直接给出单元之间的对应；需要细粒度对应时，再在这个空间上使用DTW或注意力。

表\ref{tab:vision-alignment-methods}按同一组维度比较本节的四种方法。前两行与最后一行是“学习可比较空间”的方法，中间两行是“在可比较空间中求对应”的方法。

\begin{table*}[tbp]
  \centering
  \small
  \begin{tabularx}{\linewidth}{@{}p{22mm}XXXX@{}}
    \toprule
    比较项 & 典型相关分析 & 动态时间规整 & 注意力软对齐 & 对比学习\\
    \midrule
    输出 & 两侧投影方向，整体相关系数 & 单调硬路径与代价 & 软对齐矩阵$\bm A$ & 两个编码器，共同空间\\
    对应粒度 & 全局 & 单元对单元 & 单元对单元 & 由正负例构造决定\\
    结构假设 & 线性（核与深度版本放宽） & 单调、首尾对应 & 无（单调注意力加入顺序） & 配对项比非配对项更相似\\
    监督 & 整体配对 & 无，需局部代价 & 下游目标间接训练 & 整体或局部配对\\
    防止平凡解 & 方差约束 & 不涉及 & 不涉及 & 负样本\\
    计算 & 一次广义特征分解 & $O(nm)$，带宽可降 & $O(nm)$，单调推理可降 & 每批$O(n^2d)$相似度\\
    主要失效 & 样本少时伪相关，丢弃单侧信息 & 倒序、局部代价不可靠 & 权重分散，非定位证据 & 假负例与假正例\\
    \bottomrule
  \end{tabularx}
  \caption{四种跨模态对齐方法的比较。$n,m$为两侧单元数，$d$为表示维数；对比学习一行中的$n$为批大小。各行描述方法本身的假设，不是同一数据上的实验比较。}
  \label{tab:vision-alignment-methods}
\end{table*}

回到开篇问题。对比学习可以给出画面帧与音频帧可比较的特征；在这些特征上运行带宽约束的DTW，可以把鸣响时刻对应到具体帧；字幕中“红色汽车”与画面区域的对应则可以由注意力或逐词最大匹配给出。对齐回答了“哪些证据说的是同一件事”，下一节讨论怎样把这些证据合成一个判断。

\section{多模态融合}
\label{sec:vision-alignment-fusion}

\term{多模态融合}（multimodal fusion）把来自多个模态的信息合并成一个预测。输入是各模态已经编码、并在必要的粒度上对应好的表示，例如第$t$秒的视觉特征$\bm z_{\mathrm v}\in\mathbb R^{d_{\mathrm v}}$与音频特征$\bm z_{\mathrm a}\in\mathbb R^{d_{\mathrm a}}$；输出是任务结果，例如“这一秒画面中的红车是否在鸣笛”的概率。难点在于两种信息需要以什么形式组合。有的判断只需把各自的证据加起来，例如画面和声音都表明有车；有的判断依赖两者的组合，例如“声音响起时画面中的车是否恰好停下”，单看哪一侧都无法回答。

可以用一个最小的例子说明组合方式的差别。设$z_{\mathrm a}\in\{-1,+1\}$表示录音中是否有鸣响，$z_{\mathrm v}\in\{-1,+1\}$表示画面中是否有车辆动作，目标是判断音画是否一致：两者同为$+1$或同为$-1$时$y=+1$，否则$y=-1$。任何形如$w_{\mathrm a}z_{\mathrm a}+w_{\mathrm v}z_{\mathrm v}+b$的加性打分都不能正确分类这四种情形。把$y=+1$的两种情形的打分相加，得$(w_{\mathrm a}+w_{\mathrm v}+b)+(-w_{\mathrm a}-w_{\mathrm v}+b)=2b$，正确分类要求它大于0；把$y=-1$的两种情形相加，同样得$2b$，正确分类又要求它小于0，二者矛盾。而乘积$z_{\mathrm a}z_{\mathrm v}$恰好等于$y$。这说明融合方式决定了模型能表示哪些跨模态关系，加性组合表示不了“一致性”这类交互。

图\ref{fig:vision-alignment-fusion}按合并发生的位置画出三种融合方式。早期融合在输入或浅层特征处把各模态拼在一起，交给一个网络，交互由网络自己在后续层中学习；晚期融合让各模态分别给出决策，只在最后汇总，结构简单、便于处理缺失，但交互能力最弱；中间融合在中间表示上用专门的运算显式建立交互，例如乘积、调制或注意力。三者的核心差别是“在哪里、以什么运算让模态相遇”，这直接决定参数量、交互能力和对缺失的鲁棒性。

\begin{figure*}[tbp]
  \centering
  % 三种融合位置的结构示意；宽约150mm，跨栏使用。只画模块与数据流，不表示层数。
\begin{tikzpicture}[x=1cm,y=1cm,font=\figurefont\small,
  block/.style={draw=FigureModel,fill=FigureModel!10,line width=0.8pt,
    minimum height=8mm,minimum width=20mm,align=center,rounded corners=1.8pt},
  input/.style={block,draw=FigureInput,fill=FigureInput!10,minimum width=15mm},
  output/.style={block,draw=FigureOutput,fill=FigureOutput!10,minimum width=14mm},
  arr/.style={-{Stealth[length=2mm]},BookInk,line width=0.8pt}]
  % 早期融合
  \node[anchor=west,font=\figurefont\small\bfseries] at (0,8.2) {早期融合：先拼接特征，再由一个网络预测};
  \node[input] (ea) at (1.0,7.3) {音频$\bm z_{\mathrm a}$};
  \node[input] (ev) at (1.0,6.2) {视觉$\bm z_{\mathrm v}$};
  \node[block] (ec) at (4.6,6.75) {拼接\\$[\bm z_{\mathrm a};\bm z_{\mathrm v}]$};
  \node[block] (ef) at (8.4,6.75) {共享网络$f$};
  \node[output] (ey) at (12.0,6.75) {预测$\hat y$};
  \draw[arr] (ea) -- (ec);
  \draw[arr] (ev) -- (ec);
  \draw[arr] (ec) -- (ef);
  \draw[arr] (ef) -- (ey);
  % 晚期融合
  \node[anchor=west,font=\figurefont\small\bfseries] at (0,5.1) {晚期融合：各模态先给出决策，再汇总};
  \node[input] (la) at (1.0,4.2) {音频$\bm z_{\mathrm a}$};
  \node[input] (lv) at (1.0,3.1) {视觉$\bm z_{\mathrm v}$};
  \node[block] (lfa) at (4.6,4.2) {音频网络$f_{\mathrm a}$};
  \node[block] (lfv) at (4.6,3.1) {视觉网络$f_{\mathrm v}$};
  \node[block] (lc) at (8.4,3.65) {决策汇总\\加权或相乘};
  \node[output] (ly) at (12.0,3.65) {预测$\hat y$};
  \draw[arr] (la) -- (lfa);
  \draw[arr] (lv) -- (lfv);
  \draw[arr] (lfa) -- node[above,sloped] {$p_{\mathrm a}$} (lc);
  \draw[arr] (lfv) -- node[below,sloped] {$p_{\mathrm v}$} (lc);
  \draw[arr] (lc) -- (ly);
  % 中间融合
  \node[anchor=west,font=\figurefont\small\bfseries] at (0,2.0) {中间融合：在中间表示上建立显式交互};
  \node[input] (ma) at (1.0,1.1) {音频$\bm z_{\mathrm a}$};
  \node[input] (mv) at (1.0,0.0) {视觉$\bm z_{\mathrm v}$};
  \node[block] (mi) at (4.6,0.55) {交互模块\\双线性／FiLM／交叉注意力};
  \node[block] (mg) at (8.4,0.55) {后续网络$g$};
  \node[output] (my) at (12.0,0.55) {预测$\hat y$};
  \draw[arr] (ma) -- (mi);
  \draw[arr] (mv) -- (mi);
  \draw[arr] (mi) -- (mg);
  \draw[arr] (mg) -- (my);
\end{tikzpicture}

  \caption{三种融合位置的结构示意。蓝框是各模态的输入表示，黄框是计算模块，红框是预测。早期融合在拼接后用一个共享网络，晚期融合各模态先给出决策$p_{\mathrm a},p_{\mathrm v}$再汇总，中间融合在中间表示上用显式交互运算。图中只表示数据流向，不表示层数或模块规模。}
  \label{fig:vision-alignment-fusion}
\end{figure*}

\subsection{基于特征拼接的早期融合}
\label{sec:vision-alignment-early-fusion}

\term{早期融合}（early fusion）把各模态的特征拼接成一个长向量，再交给一个网络：
\begin{equation}
  \hat y=f\bigl([\bm z_{\mathrm a};\bm z_{\mathrm v}]\bigr),
  \qquad [\bm z_{\mathrm a};\bm z_{\mathrm v}]\in\mathbb R^{d_{\mathrm a}+d_{\mathrm v}}.
  \label{eq:vision-alignment-early}
\end{equation}
它的优点是简单，而且从第一层起两种模态就在同一网络中，原则上任何交互都可以由后续层学习。

但拼接本身不产生交互。若$f$的第一层是线性的，权重$\bm W=[\bm W_{\mathrm a},\bm W_{\mathrm v}]$按列分块，则
\[
  \bm W[\bm z_{\mathrm a};\bm z_{\mathrm v}]=\bm W_{\mathrm a}\bm z_{\mathrm a}+\bm W_{\mathrm v}\bm z_{\mathrm v},
\]
两种模态在这一层只是相加。若$f$只有这一层，早期融合与对两侧分别做线性映射再相加完全等价，表示不了上面的一致性判断。交互来自之后的非线性层：例如加一层ReLU隐层，网络可以用分段线性函数逼近乘积。这说明早期融合的交互能力依赖网络深度和训练数据，没有被结构显式保证。

早期融合对输入的组织有要求。拼接要求两侧单元在同一粒度上对应：若在原始序列层面拼接，需要先把音频重采样到视频帧率，这依赖\ref{sec:vision-alignment-heterogeneity}节的时钟同步；时间偏移会让每个拼接向量中的两半描述不同时刻。各模态的数值尺度也不同，例如像素特征在$[0,1]$而音频能量的对数在$[-80,0]$，需要分别归一化，否则尺度大的模态在第一层中占主导。参数量方面，第一层的权重为$h\times(d_{\mathrm a}+d_{\mathrm v})$，与两侧维数之和成正比；以$d_{\mathrm a}=d_{\mathrm v}=512$、隐层宽$h=256$为教学设定，约26万个参数。

早期融合最大的弱点是缺失。训练时两半都有值，若推理时音频缺失而填零，网络看到的是训练中从未出现过的输入分布，输出难以预料。\ref{sec:vision-alignment-missing}节会讨论怎样在训练中为这种情形做准备。在序列模型中，把视觉词元与文字词元拼成一条序列送入同一个Transformer，是词元层面的早期融合，\ref{sec:vision-models-llava}节的LLaVA就是这样连接视觉与语言的；那里的交互由自注意力在每一层中建立，比单纯的特征拼接加多层感知机强得多。

\subsection{基于决策汇总的晚期融合}
\label{sec:vision-alignment-late-fusion}

\term{晚期融合}（late fusion）让每个模态各自训练一个预测器，最后只合并它们的输出。最简单的做法是对概率或分类分数加权平均：
\begin{equation}
  \hat p(y)=\sum_{m=1}^{M}w_m\,p_m(y\mid\bm z_m),\qquad
  w_m\geq0,\ \sum_m w_m=1.
  \label{eq:vision-alignment-late-average}
\end{equation}
权重可以在验证集上选择，也可以由一个小网络根据输入预测。晚期融合的结构很少依赖各模态的对应粒度：音频分类器可以每秒给一次结果，视频分类器可以每段给一次结果，只要最后能对到同一时间区间即可。它也天然便于处理缺失：某个模态不在，就去掉它那一项，把其余权重重新归一化。

加权平均没有概率上的依据。若假设给定类别后各模态条件独立，即$p(\bm z_{\mathrm a},\bm z_{\mathrm v}\mid y)=p(\bm z_{\mathrm a}\mid y)\,p(\bm z_{\mathrm v}\mid y)$，可以推出另一种合并方式。由贝叶斯公式，
\[
  p(y\mid\bm z_{\mathrm a},\bm z_{\mathrm v})
  =\frac{p(\bm z_{\mathrm a}\mid y)\,p(\bm z_{\mathrm v}\mid y)\,p(y)}{p(\bm z_{\mathrm a},\bm z_{\mathrm v})}.
\]
再对每个条件似然用一次贝叶斯公式，$p(\bm z_m\mid y)=p(y\mid\bm z_m)\,p(\bm z_m)/p(y)$，代入后把与$y$无关的因子并入归一化常数，得到
\begin{equation}
  p(y\mid\bm z_{\mathrm a},\bm z_{\mathrm v})\propto
  \frac{p(y\mid\bm z_{\mathrm a})\,p(y\mid\bm z_{\mathrm v})}{p(y)}.
  \label{eq:vision-alignment-poe}
\end{equation}
各模态的后验相乘，再除以一次多算的先验。二分类时取对数几率更直观：$\operatorname{logit}p(y\mid\bm z_{\mathrm a},\bm z_{\mathrm v})=\operatorname{logit}p(y\mid\bm z_{\mathrm a})+\operatorname{logit}p(y\mid\bm z_{\mathrm v})-\operatorname{logit}p(y)$，各模态的证据以对数几率相加。这种把多个分布相乘再归一化的组合称为\term{乘积专家}（product of experts），每个模态的预测器是一个“专家”\scite{vision-poe2002}它与加权平均（混合专家）的区别在于：乘积要求所有专家都同意，任何一个专家给出低概率的结果都会被压低；平均只要有一个专家支持，结果就保留一定概率。

用开篇问题的数字比较两者。设先验$p(\text{鸣笛})=0.5$，音频预测器给出$0.9$，视觉预测器看到红车刚转完弯，给出$0.4$。等权平均得$0.65$。按式\eqref{eq:vision-alignment-poe}，对数几率为$\log 9+\log(2/3)-0\approx1.792$，概率约$0.857$。音频证据很强、视觉证据略弱，乘积合并后的结论比任一单独模态更接近音频一侧，又因视觉不太支持而低于$0.9$。

当各专家是连续变量上的高斯分布时，乘积仍是高斯分布，并有闭式。设音频估计鸣笛时刻服从$\mathcal N(\mu_{\mathrm a},\sigma_{\mathrm a}^2)$，视觉估计服从$\mathcal N(\mu_{\mathrm v},\sigma_{\mathrm v}^2)$。两个密度相乘后，指数部分是$-\frac{(t-\mu_{\mathrm a})^2}{2\sigma_{\mathrm a}^2}-\frac{(t-\mu_{\mathrm v})^2}{2\sigma_{\mathrm v}^2}$；按$t$配方，$t^2$的系数为$-\frac12(\sigma_{\mathrm a}^{-2}+\sigma_{\mathrm v}^{-2})$，$t$的一次项系数为$\mu_{\mathrm a}\sigma_{\mathrm a}^{-2}+\mu_{\mathrm v}\sigma_{\mathrm v}^{-2}$，于是乘积正比于$\mathcal N(\mu,\sigma^2)$，其中
\begin{equation}
  \frac{1}{\sigma^2}=\frac{1}{\sigma_{\mathrm a}^2}+\frac{1}{\sigma_{\mathrm v}^2},\qquad
  \mu=\sigma^2\left(\frac{\mu_{\mathrm a}}{\sigma_{\mathrm a}^2}+\frac{\mu_{\mathrm v}}{\sigma_{\mathrm v}^2}\right).
  \label{eq:vision-alignment-gaussian-poe}
\end{equation}
精度（方差的倒数）相加，均值按精度加权。卷二随机变量章中高斯先验与高斯观测结合时的“按精度加权折中”是同一计算。设音频估计鸣笛在$4.3$秒、标准差$0.1$秒，视觉估计在$4.6$秒、标准差$0.3$秒，则两者精度为$100$与约$11.1$，合并后均值约$4.33$秒、标准差约$0.095$秒：结果靠近更精确的音频，且比任一单独估计更确定。Wu与Goodman的多模态变分自编码器用这一闭式合并各模态编码器给出的潜变量后验，并把先验也当作一个专家；某个模态缺失时，只需把它的专家从乘积中去掉\scite{vision-mvae2018}

乘积专家的条件独立假设在冲突时会出问题。若音频估计$4.3\pm0.1$秒，视觉估计$6.0\pm0.1$秒，两者相差17个标准差，几乎可以肯定不是同一事件，例如鸣笛来自画面外的另一辆车。式\eqref{eq:vision-alignment-gaussian-poe}却给出$5.15$秒、标准差约$0.07$秒的结论：一个两侧都不支持的时刻，而且比单侧更自信。等权混合会给出在$4.3$与$6.0$附近各有一个峰的分布，至少保留了“有两种可能”的信息。这对应\ref{sec:vision-alignment-missing-conflict}节的冲突情形：乘积合并只在各模态确实描述同一事件时可靠，使用前需要检查各专家是否相容。条件独立在另一方向上也会被违反：若两个“模态”其实来自同一信号，例如原始音频与其梅尔频谱，同一证据会被计算两次，合并结果同样过于自信。

晚期融合的参数量是各单模态预测器之和，以$d_{\mathrm a}=d_{\mathrm v}=512$、每个预测器输出256维为教学设定，两个线性头共约26万个参数，与早期融合相当。它的根本局限是交互能力：无论加权平均还是乘积，合并都发生在各自的决策之后，在对数几率上都是加性的，同样表示不了本节开头的一致性判断。两个模态需要相互参照才能理解时，必须在决策之前让它们相遇。

\subsection{基于交互建模的中间融合}
\label{sec:vision-alignment-intermediate-fusion}

\term{中间融合}（intermediate fusion）在中间表示上用专门的运算建立模态之间的交互，而不是把交互完全交给后续网络去学习。本节开头的例子表明，最直接的交互是乘积；问题在于高维向量的全部两两乘积数量巨大。下面依次讨论三类做法，它们以不同方式控制这一代价：双线性与张量融合保留所有乘积但用低秩分解压缩参数，FiLM只让一个模态调制另一个模态的每个通道，交叉注意力让一侧按内容从另一侧选择性读取。

\term{双线性池化}（bilinear pooling）对两侧向量的所有分量两两相乘，再线性组合。输出的第$k$维为
\begin{equation}
  f_k=\bm z_{\mathrm a}^\top\bm W_k\bm z_{\mathrm v}
  =\sum_{p=1}^{d_{\mathrm a}}\sum_{q=1}^{d_{\mathrm v}}W_{k,p,q}\,z_{\mathrm a,p}\,z_{\mathrm v,q},
  \qquad k=1,\ldots,K,
  \label{eq:vision-alignment-bilinear}
\end{equation}
其中$\bm W_k\in\mathbb R^{d_{\mathrm a}\times d_{\mathrm v}}$，等价地写作$f_k=\langle\bm W_k,\bm z_{\mathrm a}\bm z_{\mathrm v}^\top\rangle$，即权重与外积矩阵的内积。本节开头的一致性判断取$d_{\mathrm a}=d_{\mathrm v}=K=1$、$W=1$即可精确表示。代价是参数量$d_{\mathrm a}d_{\mathrm v}K$：以$d_{\mathrm a}=d_{\mathrm v}=512$、$K=256$为教学设定，约$6.7\times10^7$个参数，是拼接加线性层的256倍。

压缩双线性池化有两条路线。Fukui等人的多模态紧致双线性池化（multimodal compact bilinear pooling, MCB）把两个向量各自经随机的计数草图投影到高维，再用快速傅里叶变换计算两个投影的卷积，以此近似外积的随机投影，避免显式形成外积\scite{vision-mcb2016}另一条路线是低秩分解。Kim等人假设$\bm W_k=\bm U\operatorname{diag}(\bm p_k)\bm V^\top$，其中$\bm U\in\mathbb R^{d_{\mathrm a}\times d'}$、$\bm V\in\mathbb R^{d_{\mathrm v}\times d'}$由所有输出维共享，$\bm p_k\in\mathbb R^{d'}$是第$k$维的组合系数\scite{vision-mlb2017}代入式\eqref{eq:vision-alignment-bilinear}：
\[
  f_k=\bm z_{\mathrm a}^\top\bm U\operatorname{diag}(\bm p_k)\bm V^\top\bm z_{\mathrm v}
  =\sum_{r=1}^{d'}p_{k,r}\,(\bm U^\top\bm z_{\mathrm a})_r\,(\bm V^\top\bm z_{\mathrm v})_r
  =\bm p_k^\top\bigl(\bm U^\top\bm z_{\mathrm a}\circ\bm V^\top\bm z_{\mathrm v}\bigr).
\]
第二个等号把对角矩阵夹在两个向量之间的二次型展开为逐项求和，第三个等号把逐项乘积写成\term{逐元素乘积}（Hadamard product）$\circ$。把$\bm p_k$排成矩阵$\bm P\in\mathbb R^{d'\times K}$，全部输出为$\bm f=\bm P^\top(\bm U^\top\bm z_{\mathrm a}\circ\bm V^\top\bm z_{\mathrm v})$。计算顺序是：两侧各投影到$d'$维，逐元素相乘，再线性映射到输出。原论文在两侧投影后加了非线性，即$\bm f=\bm P^\top(\sigma(\bm U^\top\bm z_{\mathrm a})\circ\sigma(\bm V^\top\bm z_{\mathrm v}))+\bm b$。参数量从$d_{\mathrm a}d_{\mathrm v}K$降为$(d_{\mathrm a}+d_{\mathrm v}+K)d'$；取$d'=1024$，上例约为$1.3\times10^6$。原文指出MCB需要约16000维的输出才能减小随机投影的偏差，低秩形式则避免了这一高维输出。

双线性只含两侧分量的乘积，不含各模态单独的线性项。\term{张量融合}（tensor fusion）给每个向量末尾补一个常数1再求外积，把低阶项一起包含进来\scite{vision-tfn2017}两个模态时，
\[
  \begin{bmatrix}\bm z_{\mathrm a}\\1\end{bmatrix}
  \begin{bmatrix}\bm z_{\mathrm v}\\1\end{bmatrix}^\top
  =\begin{bmatrix}\bm z_{\mathrm a}\bm z_{\mathrm v}^\top&\bm z_{\mathrm a}\\
                  \bm z_{\mathrm v}^\top&1\end{bmatrix},
\]
右侧左上块是双线性项，右上与左下两块是两个模态各自的原始特征，右下是常数。Zadeh等人的张量融合网络对文字、视觉、声学三个模态做三重外积$[\bm z_{\mathrm l};1]\otimes[\bm z_{\mathrm v};1]\otimes[\bm z_{\mathrm a};1]$，得到的张量同时含有三个单模态块、三个双模态块和一个三模态块。原文中三个模态的表示分别为128、32、32维，张量大小为$129\times33\times33=140481$，之后接全连接层预测情感强度。若后接层的输出为64维（教学设定），仅这一层就约有$9.0\times10^6$个参数。

Liu等人的低秩多模态融合（low-rank multimodal fusion, LMF）把这一层的权重张量分解为$r$个秩一项之和，避免形成外积\scite{vision-lmf2018}记$\tilde{\bm z}_m=[\bm z_m;1]$。设输出第$k$维的权重是$\mathcal W_k=\sum_{i=1}^{r}\bm w_{\mathrm a,k}^{(i)}\otimes\bm w_{\mathrm v,k}^{(i)}$，则
\[
  h_k=\langle\mathcal W_k,\tilde{\bm z}_{\mathrm a}\otimes\tilde{\bm z}_{\mathrm v}\rangle
  =\sum_{i=1}^{r}\bigl(\bm w_{\mathrm a,k}^{(i)\top}\tilde{\bm z}_{\mathrm a}\bigr)
                 \bigl(\bm w_{\mathrm v,k}^{(i)\top}\tilde{\bm z}_{\mathrm v}\bigr).
\]
这一步用到秩一张量与外积的内积可分解为各因子内积之积。把各输出维的因子排成矩阵$\bm W_m^{(i)}\in\mathbb R^{d_h\times(d_m+1)}$，就有
\begin{equation}
  \bm h=\sum_{i=1}^{r}\bigl(\bm W_{\mathrm a}^{(i)}\tilde{\bm z}_{\mathrm a}\bigr)\circ
        \bigl(\bm W_{\mathrm v}^{(i)}\tilde{\bm z}_{\mathrm v}\bigr).
  \label{eq:vision-alignment-lmf}
\end{equation}
更多模态时在逐元素乘积中加入相应因子。参数量为$r\,d_h\sum_m(d_m+1)$，随模态数线性增长，而不是随各维数的乘积增长。沿用上面三模态的维数与64维输出，$r=4$时约为$4\times64\times(129+33+33)\approx5.0\times10^4$个参数，不到张量融合后接层的百分之一。原文在速度比较中使用$r=4$，并报告更高的秩在其数据上训练较不稳定。式\eqref{eq:vision-alignment-lmf}与前面的低秩双线性形式相同，区别在于补了常数1，因而保留了单模态项。

\term{特征逐维线性调制}（feature-wise linear modulation, FiLM）采用不对称的交互：用一个模态的表示生成参数，去调制另一个模态特征图的每个通道\scite{vision-filmcond2018}设视觉特征图第$c$个通道在位置$(i,j)$的值为$F_{c,i,j}$，条件表示为$\bm q$，例如问题文本的编码，则
\begin{equation}
  \widehat F_{c,i,j}=\gamma_c(\bm q)\,F_{c,i,j}+\beta_c(\bm q),
  \label{eq:vision-alignment-film}
\end{equation}
$\gamma_c$与$\beta_c$由一个小网络从$\bm q$预测，同一通道的所有位置共用。$\gamma_c$是乘性交互：文本可以放大、抑制甚至翻转某个视觉通道；$\beta_c$是加性偏移。原文把$\gamma$参数化为$1+\Delta\gamma$，使模块初始时接近恒等变换，因为以零为中心的$\gamma$会把激活及其梯度置零。在视觉推理数据集CLEVR上，原文的视觉部分是4个带FiLM的残差块，每块128个通道，紧邻FiLM的批归一化不再使用自己的缩放和平移参数；问题由4096维隐状态的门控循环单元读取，词嵌入为200维，末状态经仿射映射为每个残差块的$(\bm\gamma,\bm\beta)$。以两个通道为例：某位置的视觉特征为$(2,-1)$，问题“红色的物体在哪里”生成$\bm\gamma=(1.5,0)$、$\bm\beta=(0,0.5)$，调制后为$(3,0.5)$，第一个通道被放大，第二个通道的原始信号被关掉、只留偏移。FiLM的参数集中在生成器中，以$C=512$个通道、512维条件为教学设定，生成一层的$(\bm\gamma,\bm\beta)$约需$512\times1024\approx5.2\times10^5$个参数；调制本身对每个位置只做一次乘加，代价与特征图分辨率无关。它的局限是只能逐通道调制、不能按位置选择：同一通道在所有位置被同样放大，要“只看右侧的车”需要后续卷积层配合。

交叉注意力把交互做成按内容的选择性读取：一侧的每个单元按与另一侧各单元的相关程度加权读取，权重就是\ref{sec:vision-alignment-attention}节的软对齐矩阵。与前两种做法相比，它的交互随位置变化，能表达“这个词对应那个区域”；而且两侧单元数可以不同，不需要预先对齐到同一粒度。单头交叉注意力的参数是查询、键、值与输出四个投影，宽度$d=512$时约$4d^2\approx1.0\times10^6$；计算量与两侧单元数之积成正比，为$O(mnd)$，这是它与前两种做法最大的差别：双线性与FiLM的代价与序列长度无关或线性相关，交叉注意力在长视频、高分辨率图像上会迅速变贵。BLIP-2的Q-Former与Flamingo的门控交叉注意力都是以交叉注意力实现的中间融合，后者还用一个初始为零的门控，使新加入的视觉读取在训练开始时不扰动冻结的语言模型（\ref{sec:vision-models-flamingo}节）。MulT则在多种模态之间成对使用交叉注意力。

表\ref{tab:vision-alignment-fusion-compare}在同一组教学维数下比较本节讨论的融合方式。比较的维度是参数量、交互形式、对输入对齐的要求以及缺失时的处理。

\begin{table*}[tbp]
  \centering
  \small
  \begin{tabularx}{\linewidth}{@{}p{24mm}p{20mm}XXX@{}}
    \toprule
    融合方式 & 参数量（教学） & 交互形式 & 对齐要求 & 模态缺失时\\
    \midrule
    拼接加线性层 & $2.6\times10^5$ & 加性；交互需后续非线性层 & 同一粒度、同一时刻 & 填充值偏离训练分布\\
    晚期融合（两个线性头） & $2.6\times10^5$ & 对数几率上加性 & 只需结果落在同一区间 & 去掉该项并重新归一化\\
    完全双线性 & $6.7\times10^7$ & 全部两两乘积 & 同一粒度 & 乘积项整体失效\\
    低秩双线性 & $1.3\times10^6$ & 秩$d'$内的两两乘积 & 同一粒度 & 同上\\
    FiLM & $5.2\times10^5$ & 条件模态逐通道乘加 & 条件为整体向量即可 & 条件缺失时可退回$\gamma=1,\beta=0$\\
    单头交叉注意力 & $1.0\times10^6$ & 按内容逐单元读取 & 不要求，单元数可不同 & 被读取侧缺失则无内容可读\\
    \bottomrule
  \end{tabularx}
  \caption{融合方式的教学比较。设两侧特征均为512维，输出或隐层为256维；低秩双线性取$d'=1024$；FiLM按512个通道、512维条件计一层生成器；交叉注意力按单头、宽度512计四个投影。参数量忽略偏置，只用于比较数量级。}
  \label{tab:vision-alignment-fusion-compare}
\end{table*}

表中可以读出一条主线：从拼接、晚期汇总到双线性，交互能力逐步增强，参数也急剧增加；低秩分解与FiLM用结构假设把参数降回可用范围；交叉注意力把代价从参数转移到了计算上。选择时应先问任务需要什么交互。只需合并独立证据时，晚期融合足够，而且在缺失时最稳健；需要两模态相互参照、但都已汇总成整体向量时，低秩双线性或FiLM更省；需要细粒度的“哪个词对应哪块区域”时，交叉注意力几乎是唯一选择。\ref{chap:vision-multimodal-models}章的模型多数采用交叉注意力或共享序列，\ref{chap:vision-multimodal-generation}章的条件生成大量使用FiLM式的调制与交叉注意力，把文字或时间条件接入生成主干。

\section{模态缺失与不平衡下的学习}
\label{sec:vision-alignment-missing}

前两节默认每个模态都在场，而且训练会充分利用每个模态。本节放松这两个条件。第一种情形是训练数据中部分样本缺少某个模态，需要决定怎样用这些不完整样本，以及怎样让模型为缺失做好准备。第二种情形是推理时某个模态缺失，而模型是在完整数据上训练的，需要在不重新训练或少量训练的前提下给出可靠预测。第三种情形是所有模态都在，但联合训练时一个模态学得快、另一个学得慢，最终模型几乎只依赖前者。三种情形的条件不同，可用的做法也不同：前两种关心输入不完整，第三种关心优化过程本身。

\subsection{训练时模态缺失下的学习}
\label{sec:vision-alignment-missing-training}

设训练集中每个样本带有掩码$\bm r\in\{0,1\}^M$。最简单的做法是只用完整样本，即$\bm r=\bm 1$的样本。这丢弃了数据；若缺失属于\ref{sec:vision-alignment-missing-conflict}节说的非随机缺失，例如低质量视频更常缺字幕，留下的完整样本还会偏向某些情形，训练出的模型在被丢弃的那类数据上表现未知。

\term{模态丢弃}（modality dropout）反过来，在完整样本上人为制造缺失。Neverova等人提出的ModDrop在训练手势识别网络时，以一定概率把整个模态通道的输入置零，每个模态独立地决定保留或丢弃；丢弃作用在各模态专属路径的输入上，融合层之前\scite{vision-moddrop2016}这与卷五神经网络的训练章中的Dropout同属随机计算，区别在于丢弃的单位是整个模态而不是单个神经元，可参考\BookChapterRef{chap:neural-network-training}{卷五“神经网络的训练”章}。从期望损失看，设各模态以概率$q$独立丢弃，丢弃后的输入为$\bm r\odot\bm z$（按模态整体置零），训练目标为
\begin{equation}
  \mathbb E_{\bm r}\bigl[L\bigl(f(\bm r\odot\bm z),y\bigr)\bigr]
  =\sum_{\bm r\in\{0,1\}^M}q^{M-\lvert\bm r\rvert}(1-q)^{\lvert\bm r\rvert}\,
   L\bigl(f(\bm r\odot\bm z),y\bigr),
  \label{eq:vision-alignment-moddrop}
\end{equation}
$\lvert\bm r\rvert$是保留的模态数。它是所有缺失模式下损失的加权和：同一组参数被要求在每种模态子集上都能预测。两个模态、$q=0.1$时，完整输入占权重$0.81$，只有一种模态的两种情形各占$0.09$，两个都缺的情形占$0.01$；实践中常把全部丢弃的情形排除。原文报告这样训练的网络在某个模态缺失或受损时性能下降更小，并认为它还促使网络学习跨模态的相关，抑制各模态专属表示之间不必要的协同适应。丢弃概率在验证集上选择：太大会让完整输入下的训练不足，太小则缺失情形的权重太低。

模态丢弃让模型学会“没有某个模态也能预测”，但没有利用缺失样本本身。另一条路线是插补与重建：用已观测的模态估计缺失模态的表示，再按完整样本处理。简单的插补用训练集上该模态特征的均值；更强的插补训练一个从模态$\mathrm a$预测模态$\mathrm v$表示的网络，在完整样本上监督它，在缺失样本上用它的输出补全。这本质上是跨模态转换。它的风险是插补值被当作真实观测：若视觉表示是由音频推出来的，融合模型再从中读出“视觉证据”，得到的只是音频信息的重复，会造成前面讨论的乘积专家式的过度自信。

\ref{sec:vision-alignment-late-fusion}节的多模态变分自编码器提供了一种不需要插补的训练方式。它的联合后验是各模态专家与先验的乘积，缺失的模态直接从乘积中去掉，于是任意模态子集都有定义好的后验。训练时，即使样本是完整的，也只在每一步计算部分证据下界项：使用全部模态的一项、每个单模态各一项，以及随机选取的若干模态子集各一项。这样每个单模态编码器都被单独训练过，不完整样本也可以按它实际拥有的模态加入训练。代价是每步要计算多个下界项，以及条件独立假设在模态冲突时的问题。

\subsection{推理时模态缺失下的预测}
\label{sec:vision-alignment-missing-inference}

训练数据完整、推理时缺失，是部署中常见的情形：摄像头被遮挡，麦克风关闭，字幕服务不可用。此时能做什么，取决于训练时选择的融合方式。

晚期融合与乘积专家的处理最直接：去掉缺失模态的预测，其余项按原规则合并。加权平均时，设原权重为$w_m$、可用模态集合为$S$，新权重为$w_m/\sum_{l\in S}w_l$。这一做法隐含的假设是，各模态的预测器在单独使用时仍然校准；若视觉预测器在训练中总有音频陪伴，它的输出未必是可靠的独立概率。早期融合和中间融合则没有这样的天然退路：缺失的那半输入必须填入某个值，而训练中未见过的填充值会让网络处于分布之外。若训练时用了模态丢弃，置零就是网络见过的情形，这是两节之间的联系。

当模型很大、无法为缺失情形重新训练时，可以只训练少量附加参数。Lee等人在冻结的视觉语言Transformer ViLT上研究这一问题\scite{vision-missingprompt2023}缺失的文字用空字符串代替，缺失的图像用像素值全部置为1的空白图代替，这样输入格式不变。对图文两种模态，共有完整、缺文字、缺图像三种情形，作者为每种情形各学习一组\term{缺失感知提示}（missing-aware prompt），按样本的实际缺失情形选择其中一组。提示是附加在Transformer层上的可学习向量：一种做法把它们拼在每层输入序列的前面，另一种拆成键和值的附加项放进注意力中，保持序列长度不变；默认加在前6层，长度为16。训练时只更新提示、汇聚层和分类头，提示约占ViLT参数的0.2\%。这种做法与\ref{sec:vision-alignment-missing-training}节的模态丢弃目标相同，都让模型适应不同缺失模式，区别在于把适应集中在少量参数上，主干不动。

视觉语言模型中还有一种隐蔽的缺失：图像在输入中，却没有被有效使用，或者图像缺失时模型仍然流畅作答。语言模型可以只靠问题中的常见搭配给出看似合理的答案，例如问“草是什么颜色”，不看图也回答“绿色”。这种情形不会报错，只能用\ref{sec:vision-understanding-input-protocol}节的移除、替换实验发现。\ref{sec:vision-models-modality-contribution}节讨论了模型选型时的相应检查。从学习的角度看，它是下一小节模态不平衡的一种表现。

\subsection{模态不平衡下的优化}
\label{sec:vision-alignment-imbalance}

直觉上，多一个模态至少不应该更差：网络可以忽略没用的输入。实际训练却常常相反。Wang等人在多个任务上观察到，端到端训练的多模态网络常常不如其中最好的单模态网络；在Kinetics动作识别上，只用画面的网络top-1准确率为72.6\%，画面加音频的朴素晚期融合网络为71.4\%\scite{vision-gblend2020}他们把原因归结为过拟合：多模态网络容量更大，训练准确率更高而验证准确率更低；而且不同模态过拟合和泛化的速度不同，用同一套优化设置联合训练对每个模态都不是最合适的。我们把联合训练中一个模态主导学习、其他模态学得不足的现象称为\term{模态不平衡}（modality imbalance）。

不平衡可以从梯度上推导出来。设一个分类网络把两侧编码器的输出拼接后接线性层，分类分数为
\[
  \bm o=\bm W_{\mathrm a}\bm\phi_{\mathrm a}+\bm W_{\mathrm v}\bm\phi_{\mathrm v}+\bm b,
\]
其中$\bm\phi_{\mathrm a}=\bm\phi_{\mathrm a}(\bm x_{\mathrm a};\bm\theta_{\mathrm a})$、$\bm\phi_{\mathrm v}$是两侧编码器的输出，$\bm W_{\mathrm a}$、$\bm W_{\mathrm v}$是线性层权重按列分出的两块。由\ref{sec:vision-alignment-early-fusion}节，拼接后的线性层正是这一加性形式。softmax交叉熵对$\bm o$的梯度为$\bm p-\bm e_y$，$\bm p$是预测概率，$\bm e_y$是真实类别的独热向量。由链式法则，
\begin{equation}
  \frac{\partial L}{\partial\bm W_{\mathrm a}}=(\bm p-\bm e_y)\bm\phi_{\mathrm a}^\top,\qquad
  \frac{\partial L}{\partial\bm\phi_{\mathrm a}}=\bm W_{\mathrm a}^\top(\bm p-\bm e_y),
  \label{eq:vision-alignment-imbalance-grad}
\end{equation}
视觉一侧同理。关键在于两侧共用同一个残差$\bm p-\bm e_y$，而它由两侧之和决定。若视觉一侧学得快，很快使$\bm p$接近$\bm e_y$，残差变小，音频一侧收到的梯度也随之变小，即使音频编码器自己还远没有学好。

\begin{example}[主导模态压低弱模态的梯度]
  \label{ex:vision-alignment-imbalance}
  以下为教学构造，用二分类与logistic输出简化。设某样本$y=1$，视觉分支已贡献分数$w_{\mathrm v}\phi_{\mathrm v}=3$，音频分支贡献$w_{\mathrm a}\phi_{\mathrm a}=0.5$，其中$w_{\mathrm a}=1$、$\phi_{\mathrm a}=0.5$。联合分数为3.5，$p=\sigma(3.5)\approx0.971$，音频权重的梯度为$(p-1)\phi_{\mathrm a}\approx-0.0147$。若音频分支单独训练，分数为0.5，$p=\sigma(0.5)\approx0.622$，梯度为$(p-1)\phi_{\mathrm a}\approx-0.189$，约为联合训练时的13倍。同样的音频特征，在联合训练中得到的学习信号只有单独训练的约十三分之一。
\end{example}

这个例子说明，不平衡不需要任何模态“没用”。只要一个模态更容易拟合训练数据，它就会先把残差降下来，其他模态便学不到它们本来能学到的东西。推理时若主导模态缺失或受干扰，弱模态无法接替。\ref{sec:vision-alignment-missing-inference}节的视觉问答例子是同一现象：语言先验很容易拟合，视觉证据就学得不足。Goyal等人为此构造了VQA v2，为同一问题配上答案不同的相似图像，使只看问题无法答对\scite{vision-vqav22017}这是从数据上纠正不平衡；下面两种方法则从优化上纠正。

Wang等人提出的梯度混合（gradient blending, G-Blend）给每个单模态分支和融合分支各接一个分类头，总损失为各头损失的加权和$L=\sum_i w_iL_i$。权重按各分支在训练中的表现确定：在一段训练中，记某分支验证损失的下降为泛化收益$G_i$，训练与验证损失之差的增加为过拟合量$O_i$，取$w_i\propto G_i/O_i^2$。泛化收益大、过拟合小的分支得到更大的权重。权重可以先用一次试训练估计后固定，也可以每隔若干轮重新估计。单模态头的损失直接监督每个编码器，使它们不再只靠共享残差获得梯度。

Peng等人的在线梯度调制（on-the-fly gradient modulation, OGM-GE）不加额外的头，而是直接减慢主导模态的更新\scite{vision-ogmge2022}它利用拼接线性层的分块结构，为每个样本计算各模态单独的置信度：把线性层中属于该模态的那一块单独作用在它的特征上，偏置平分，取softmax后在真实类别上的概率，
\[
  s_i^{\mathrm a}=\operatorname{softmax}\bigl(\bm W_{\mathrm a}\bm\phi_{\mathrm a}(\bm x_i^{\mathrm a})+\bm b/2\bigr)_{y_i},
\]
视觉一侧$s_i^{\mathrm v}$同理。在一个小批$B_t$上，两模态的差异比为
\[
  \rho_t^{\mathrm v}=\frac{\sum_{i\in B_t}s_i^{\mathrm v}}{\sum_{i\in B_t}s_i^{\mathrm a}},\qquad
  \rho_t^{\mathrm a}=\frac{1}{\rho_t^{\mathrm v}}.
\]
对每个模态$u\in\{\mathrm a,\mathrm v\}$，调制系数与编码器参数的更新为
\begin{equation}
  k_t^{u}=\begin{cases}1-\tanh(\alpha\rho_t^{u}),&\rho_t^{u}>1,\\1,&\text{否则},\end{cases}
  \qquad
  \bm\theta_{t+1}^{u}=\bm\theta_t^{u}-\eta\,k_t^{u}\,\tilde{\bm g}(\bm\theta_t^{u}),
  \label{eq:vision-alignment-ogm}
\end{equation}
其中$\tilde{\bm g}$是小批随机梯度，$\eta$是学习率，$\alpha>0$是控制调制强度的超参数，原文在验证集上从$[0,1]$中选取。置信度较高的模态$\rho>1$，其更新被乘以小于1的系数而放慢；较弱的模态保持原速，于是在共享残差尚未降到零之前有更多时间学习。设一个小批两个样本的视觉置信度为$(0.9,0.8)$、音频置信度为$(0.5,0.3)$，则$\rho^{\mathrm v}=1.7/0.8=2.125$，$\alpha=0.5$时$k^{\mathrm v}=1-\tanh(1.0625)\approx0.213$，视觉编码器本步只走原步长的约五分之一；$\rho^{\mathrm a}\approx0.47<1$，音频编码器不受影响。注意式\eqref{eq:vision-alignment-ogm}在$\rho=1$处不连续：$\rho$略大于1时系数就跳到$1-\tanh\alpha$。

减小梯度的幅度也减小了随机梯度噪声的幅度，原文认为这会削弱随机梯度噪声带来的泛化作用。GE（generalization enhancement）一项因此在调制后的梯度上加一个零均值高斯噪声$\bm h\sim\mathcal N(\bm 0,\bm\Sigma^{\mathrm{sgd}})$，其协方差取为随机梯度噪声的协方差，每步重新抽取；更新变为$\bm\theta_{t+1}^{u}=\bm\theta_t^{u}-\eta(k_t^{u}\tilde{\bm g}+\bm h)$。原文在CREMA-D、Kinetics-Sounds、VGGSound等音画数据集上，以ResNet18为两侧编码器（音频分支的输入通道由3改为1），在拼接、求和与FiLM等融合方式上验证了这一方法。

两种方法的适用条件不同。G-Blend需要额外的单模态头和估计权重的试训练，但不依赖融合层的具体形式。OGM-GE不增加参数，但它的置信度计算依赖“融合层可以按模态分块”这一结构，在拼接、求和这类加性融合上最自然；对于交叉注意力这类深度交互的融合，单模态置信度没有现成的定义。两者都把“某模态置信度高”当作“该模态学得太快”的信号，若一个模态本来就更有信息量，压低它的学习未必有利，因此调制强度需要在验证集上确定。

\section{本章小结}
\label{sec:vision-alignment-summary}

开篇的问题是：在画面、录音和字幕三种记录中，判断鸣笛的是不是那辆右转的红车。本章把这个问题拆成对齐、融合与缺失不平衡三部分，分别给出了回答。

对齐回答“哪些单元说的是同一件事”。多模态数据在数值形式和时间分辨率上异构，配对往往只在整体层面成立且带有噪声，所以对应关系需要建立而不能直接读出。四种方法分两类：典型相关分析和对比学习学习一个可比较的空间，前者是线性投影上的广义特征值问题，后者用负样本防止塌缩，其损失给出互信息的下界并受$\log K$限制；动态时间规整和注意力在可比较的空间中求对应，前者以动态规划求单调的硬路径，后者以softmax权重给出可微的软对应，单调注意力把二者的长处结合起来。融合回答“对上的证据怎样合成判断”。早期融合拼接特征，加性的第一层要靠后续非线性层才能学到交互；晚期融合汇总决策，乘积专家在条件独立时有概率依据，冲突时却会过于自信；中间融合用双线性、张量、FiLM或交叉注意力显式建立交互，低秩分解把参数降到可用范围，交叉注意力则把代价转到计算上。缺失与不平衡回答“模态不全或学得不均时怎么办”：模态丢弃与乘积专家让模型在训练中见到各种缺失模式，缺失感知提示以少量参数适应推理时的缺失，而联合训练中共用的残差会让主导模态压低弱模态的梯度，G-Blend与OGM-GE分别用额外监督和梯度调制纠正这一点。

对开篇问题，这些方法能做到的是：用对比学习得到的音画特征加上DTW，把鸣响对到具体帧；用注意力把“红色汽车”对到画面区域；用中间融合判断音画是否一致，并在录音缺失时退回视觉判断。它们没有解决的部分同样重要。对齐方法只说明哪些单元相似或共现，鸣笛声是否“来自”那辆车仍需要声源定位这样的空间证据；融合在冲突时没有通用的裁决规则，常常需要显式报告不一致；本章的方法都依赖训练数据覆盖的配对与缺失模式，训练中未出现的模态组合仍需单独检验。CLIP、BLIP-2、Flamingo等模型怎样把这些原理落成具体结构，见\ref{chap:vision-multimodal-models}章；在检索、定位与问答中怎样使用并核验对齐结果，见\ref{chap:vision-multimodal-understanding}章；音画同步生成怎样利用时间对齐的条件，见\ref{chap:vision-multimodal-generation}章。

\paragraph{扩展阅读}
系统梳理方面，Liang等人的综述在本章所依据的Baltrušaitis等人的分类之后，把多模态学习的核心挑战重新组织为表示、对齐、推理、生成、迁移与量化六类，并讨论了模态异构、模态间联系与交互的含义；Xu等人的综述专门整理以Transformer为骨架的多模态模型，从输入的词元化、融合在网络中的位置到预训练任务逐项比较；Xu等人更早的多视图学习综述按协同训练、多核学习与子空间学习三条路线，说明同一对象的多个视图怎样相互提供监督，是理解CCA一类方法的另一种背景\scite{vision-liang2024foundations,vision-xu2023mmtransformer,vision-xu2013multiview}

对齐方向，Wang等人在深度CCA之外比较了把相关性目标与自编码器重构结合的若干目标，并讨论了用小批量随机梯度优化相关性目标的做法\scite{vision-wang2015dccae}需要在海量数据中检索时，可以读Jiang与Li的深度跨模态哈希，它端到端地把图像与文本映射到共享的二值码，使跨模态检索可以用汉明距离完成\scite{vision-jiang2017dcmh}DTW的硬路径不可微，Cuturi与Blondel的soft-DTW把动态规划递推中的最小值换成带温度的软最小值，得到可微的对齐损失，可以接在网络之后训练；Graves等人的连接时序分类（connectionist temporal classification, CTC）则在只知道输出序列、不知道逐帧对应时，对所有单调对齐路径求和来定义似然，是语音和手写识别中处理未分段序列的常用做法\scite{vision-cuturi2017softdtw,nlp-speech-graves2006}对没有顺序的区域与词，Chen等人的图最优传输把对齐写成两组单元之间的运输方案，同时匹配单元本身与单元之间的关系。它可以与\ref{sec:vision-alignment-attention}节的软对齐对照：运输方案受两侧边缘分布的约束，每一侧的质量都要分配完，而注意力权重只在源一侧归一化\scite{vision-chen2020got}

融合方向，Arevalo等人的门控多模态单元由输入计算门控值，再按门控值对各模态的表示加权求和；与\ref{sec:vision-alignment-early-fusion}节的拼接相比，它让每个样本按自身内容决定各模态所占的比重。Nagrani等人的注意力瓶颈只在Transformer的后几层融合音频与视觉，并且只允许两者通过少量共享的瓶颈词元交换信息，既保留了中间融合的交互，又把跨模态注意力的计算量控制在较小范围\scite{vision-arevalo2017gmu,vision-nagrani2021mbt}
